% Options for packages loaded elsewhere
\PassOptionsToPackage{unicode}{hyperref}
\PassOptionsToPackage{hyphens}{url}
\documentclass[
  twocolumn]{article}
\usepackage{xcolor}
\usepackage[top=1in,bottom=1in,left=0.75in,right=0.75in]{geometry}
\usepackage{amsmath,amssymb}
\setcounter{secnumdepth}{5}
\usepackage{iftex}
\ifPDFTeX
  \usepackage[T1]{fontenc}
  \usepackage[utf8]{inputenc}
  \usepackage{textcomp} % provide euro and other symbols
\else % if luatex or xetex
  \usepackage{unicode-math} % this also loads fontspec
  \defaultfontfeatures{Scale=MatchLowercase}
  \defaultfontfeatures[\rmfamily]{Ligatures=TeX,Scale=1}
\fi
\usepackage{lmodern}
\ifPDFTeX\else
  % xetex/luatex font selection
\fi
% Use upquote if available, for straight quotes in verbatim environments
\IfFileExists{upquote.sty}{\usepackage{upquote}}{}
\IfFileExists{microtype.sty}{% use microtype if available
  \usepackage[]{microtype}
  \UseMicrotypeSet[protrusion]{basicmath} % disable protrusion for tt fonts
}{}
\makeatletter
\@ifundefined{KOMAClassName}{% if non-KOMA class
  \IfFileExists{parskip.sty}{%
    \usepackage{parskip}
  }{% else
    \setlength{\parindent}{0pt}
    \setlength{\parskip}{6pt plus 2pt minus 1pt}}
}{% if KOMA class
  \KOMAoptions{parskip=half}}
\makeatother
\setlength{\emergencystretch}{3em} % prevent overfull lines
\providecommand{\tightlist}{%
  \setlength{\itemsep}{0pt}\setlength{\parskip}{0pt}}
\AtBeginDocument{\fontsize{9}{10.5}\selectfont}
\usepackage{graphicx}
\usepackage{booktabs}
\usepackage{array}
\usepackage{float}
\usepackage{xcolor}
\usepackage{caption}
\usepackage{bookmark}
\IfFileExists{xurl.sty}{\usepackage{xurl}}{} % add URL line breaks if available
\urlstyle{same}
\hypersetup{
  pdftitle={Spatial Configuration and Social Behavior in LLM-Agent Systems},
  pdfauthor={Anonymous Authors},
  hidelinks,
  pdfcreator={LaTeX via pandoc}}

\title{Spatial Configuration and Social Behavior in LLM-Agent Systems}
\usepackage{etoolbox}
\makeatletter
\providecommand{\subtitle}[1]{% add subtitle to \maketitle
  \apptocmd{\@title}{\par {\large #1 \par}}{}{}
}
\makeatother
\subtitle{Compact Conference-Style Draft}
\author{Anonymous Authors}
\date{2026-05-01}

\begin{document}
\maketitle

\section{Introduction}\label{introduction}

\begin{table*}[t]
\centering
\caption{Table 1. Multi-Survey Positioning Matrix}
\scriptsize
\setlength{\tabcolsep}{3pt}
\renewcommand{\arraystretch}{1.12}
\begin{tabular}{@{}>{\raggedright\arraybackslash}p{0.175\textwidth}|>{\raggedright\arraybackslash}p{0.175\textwidth}|>{\raggedright\arraybackslash}p{0.175\textwidth}|>{\raggedright\arraybackslash}p{0.175\textwidth}|>{\raggedright\arraybackslash}p{0.175\textwidth}@{}}
\toprule
Neighboring review area & Representative local source(s) & Main focus in that review space & How space is usually treated & Gap this survey targets \\
\midrule
Spatial intelligence across scales & 01\_Feng2025\_Spatial\_Intelligence\_Across\_Scales.pdf & Broad spatial intelligence, smart-city contexts, embodied and multi-scale AI systems & Space is a broad capability domain across scales and applications & Does not center agent-accessible spatial representation in LLM-agent social simulation. \\
LLM multi-agent systems & 02\_Guo2024\_LLM\_Multi\_Agents\_Survey.pdf & Agent architectures, collaboration, communication, planning, and multi-agent coordination & Space is one possible environment feature among many & Does not provide a fine-grained spatial representation taxonomy or evidence map. \\
LLM game agents and NPCs & 03\_Hu2024\_LLM\_Game\_Agents\_Survey.pdf & Game agents, NPC behavior, sandbox worlds, and perception-action loops & Space appears through game worlds, locations, and embodied interaction & Does not separate backend world richness from what the agent actually receives. \\
LLM-ABM and simulation & 04\_Gao2024\_LLM\_ABM\_Simulation\_Survey.pdf & Connections between LLM agents, ABM, simulation domains, and evaluation challenges & Environment structure is part of simulation design & Does not isolate configuration-level or agent-facing spatial input as the main analytic object. \\
LLM social simulation & 05\_Mou2024\_Social\_Simulation\_LLM\_Agents\_Survey.pdf & Individual, scenario, and society-level social simulation with LLM agents & Social context is central; spatial structure is usually secondary & Does not ask whether spatial representation mediates social behavior. \\
Social-agent evaluation and game theory & 06\_Feng2024\_Social\_Agents\_Game\_Theory\_Survey.pdf & Evaluation, interaction settings, and game-theoretic social behavior & Space is not the primary organizing dimension & Useful for evaluation framing, but not for spatial representation coverage. \\
General LLM-agent methodology & 07\_Luo2025\_LLM\_Agent\_Methodology\_Survey.pdf & Agent methodology, collaboration, evolution, and evaluation ecosystems & Environment is part of agent design, not the central variable & Does not map L0-L5 agent-facing spatial representation. \\
VLN and navigation & 08\_Gu2022\_VLN\_Survey.pdf; 09\_Zhang2024\_VLN\_Foundation\_Models\_Survey.pdf & Navigation tasks, embodied perception, route following, and foundation-model era navigation & Space is central, but mostly as navigation and world-modeling & Navigation success does not directly establish spatially mediated social behavior. \\
VLA / embodied AI & 10\_Ma2024\_VLA\_Embodied\_AI\_Survey.pdf & Vision-language-action systems, embodied interfaces, robotic or simulated action & Geometry and embodiment are often central & Embodied action evidence does not by itself answer LLM social-simulation claims. \\
LLM-agent scoping / architecture reviews & 11\_Silacci2026\_LLM\_Agents\_Scoping\_Review.pdf; 12\_Leiser2025\_LLM\_Architectures\_Scoping\_Review.pdf; 13\_TudorCar2020\_Conversational\_Agents\_Scoping\_Review.pdf & Scoping-review method, architecture dimensions, conceptual mapping, and outcome tables & Space is usually not the focal taxonomy & Provides reporting models, but not the WHERE-focused evidence map. \\
This survey & Current evidence map and appendix assets & Agent-accessible spatial representation in LLM-agent social simulation & Space is coded by what agents receive: L0-L5, backend vs interface, layer and evidence status & Defines the WHERE gap: current systems stage social worlds, but rarely expose agent-facing global configuration or evaluate it under controlled social-behavior comparisons. \\
\bottomrule
\end{tabular}

\label{tab:table-1-multi-survey-positioning-matrix}
\end{table*}

\begin{figure*}[t]
\centering
\includegraphics[width=0.94\textwidth,height=0.62\textheight,keepaspectratio]{../figures/gpt\_image\_2/figure\_1\_where\_gap\_claim\_architecture\_gpt\_image\_2\_v3.png}
\caption{Figure 1. WHERE gap and evidence-role architecture. The survey maps a missing intersection between spatial-configuration theory, LLM spatial capability work, and LLM multi-agent social simulation: agent-facing spatial configuration linked to observed social behavior. The main evidence map is built from `Core` rows, separated into strict `anchor\_core` and lower-weight `bridge\_core` layers. `Adjacent` work supports feasibility and boundary reasoning, while `Foundational` work supports theory and transferable hypotheses. The figure defines a descriptive and gap-oriented claim architecture rather than a causal proof.}
\label{fig:1}
\end{figure*}

LLM-based agents are increasingly used to build simulated towns, game
worlds, virtual communities, social NPCs, embodied assistants, and
agent-based social simulations. Many of these systems are spatial in an
everyday sense: agents occupy places, move through environments,
encounter others nearby, respond to scene context, or act inside
rendered worlds. Yet the presence of a spatial world does not by itself
answer a more specific question: what spatial structure is actually
available to the agent, and what evidence shows that this structure
changes social behavior?

This review focuses on that question. Its starting point is a gap
between three neighboring literatures. Space Syntax and related
spatial-simulation traditions provide concepts for studying how
configuration, accessibility, depth, visibility, and route structure may
shape movement and encounter in physical environments. LLM
spatial-reasoning and spatially aware agent work asks whether models can
interpret or act on spatial inputs. LLM multi-agent and
social-simulation work builds populations of language agents that
interact, coordinate, move, and sometimes produce emergent social
patterns. What remains underdeveloped is the intersection: agent-facing
spatial configuration linked to observed social behavior in LLM
multi-agent systems.

We call this the WHERE gap. The issue is not simply whether LLM agents
can answer spatial questions, nor whether multi-agent systems can be
placed inside maps or 3D environments. The issue is where spatial
information enters the agent loop: as a place label, a semantic scene
description, local adjacency, co-presence, global abstract structure,
geometry-bearing perception, or something else. Without this
interface-level distinction, a visually rich simulator can be mistaken
for a configurationally rich agent system, and a reported social pattern
can be overread as evidence of spatial mediation.

\subsection{1.1 Research Question and
Scope}\label{research-question-and-scope}

The guiding question is:

\begin{quote}
What is known, what is missing, and what is needed to study how spatial
configuration may shape social behavior in LLM multi-agent systems?
\end{quote}

The wording is deliberately cautious. This paper does not ask whether
spatial configuration has already been shown to shape LLM-agent
societies in a general or causal sense. The current evidence base is not
mature enough for that claim. Instead, the paper asks a scoping-review
question: how existing systems represent space at the agent interface,
which social behaviors they study, what evidence status their results
support, and what future studies would need before stronger
spatial-behavior claims become available.

The review therefore treats ``space'' as an agent-accessible
representation problem rather than only an environment-design problem. A
system may contain a 3D engine, a graph backend, a grid, a map, or a
narrative setting. Those backend forms matter, but they are not the
coding target by themselves. The coding target is what the agent can
consume while perceiving, reasoning, planning, communicating, or acting.
This is why the review distinguishes environment-side richness from
agent-facing spatial structure throughout the paper.

The social side is also bounded. The review is concerned with socially
relevant behavior: dialogue, co-presence, cooperation, conflict, partner
choice, mobility with social consequences, group formation, role
differentiation, information exposure, or emergent social structure.
Pure navigation, path planning, robot control, or single-agent spatial
reasoning is included only as adjacent feasibility evidence unless it
connects directly to multi-agent social behavior.

\subsection{1.2 Relation to Existing
Surveys}\label{relation-to-existing-surveys}

Existing surveys cover much of the surrounding territory.
Spatial-intelligence surveys map broad capability domains across scales.
LLM multi-agent surveys emphasize agent architectures, collaboration,
communication, planning, and coordination. Game-agent and NPC surveys
discuss sandbox worlds, perception-action loops, and character behavior.
LLM-ABM and social-simulation surveys organize the emerging simulation
literature around individuals, scenarios, societies, evaluation, and
applications. Navigation and embodied-AI surveys treat space as central,
but usually through route following, perception, action, and task
success rather than social behavior.

Table 1 positions this review against those neighboring survey areas.
The difference is not that previous surveys ignore environments
entirely. Many of them discuss environments, embodiment, games,
navigation, and simulation settings. The difference is the analytic
object. This review centers agent-accessible spatial representation as
the bridge between environment design and social-behavior claims. It
asks what agents receive, what behavior is measured, and what level of
evidence connects the two.

This positioning also clarifies what the paper does not try to do. It is
not a general survey of LLM agents, a full survey of spatial reasoning
benchmarks, a comprehensive Space Syntax tutorial, or a navigation
survey. It uses each of those literatures only to the extent needed for
the WHERE problem. Space Syntax supplies transferable concepts and
hypotheses. Adjacent spatial-reasoning work supplies feasibility
boundaries. Core LLM-agent social-simulation papers supply the evidence
map.

\subsection{1.3 Review Design and Evidence
Roles}\label{review-design-and-evidence-roles}

The review is designed as a scoping review with a structured evidence
map rather than as an effect-size synthesis. The corpus is separated
into three evidence roles before claims are made. \texttt{Core} papers
describe LLM multi-agent systems with identifiable spatial environments
and socially relevant behavior. These form the main evidence map.
\texttt{Adjacent} papers include spatial reasoning, spatially aware
agents, embodied or boundary systems, and related work that informs
feasibility or scope. They do not count as direct evidence for
multi-agent social effects. \texttt{Foundational} papers include Space
Syntax, physical-space empirical work, and earlier spatial-simulation
traditions. They motivate hypotheses and vocabulary, not direct claims
about LLM-agent systems.

Figure 1 summarizes this role separation. Within the main evidence map,
the Core corpus is further divided into a strict \texttt{anchor\_core}
and a widened \texttt{bridge\_core}. The strict anchor supports the
cleanest descriptive mapping and strict-gap claims. The bridge layer
broadens coverage but carries lower evidential weight and must remain
explicitly qualified. Boundary cases such as \texttt{HC01} and
\texttt{TW-02} are retained for scope discipline but do not enter the
stable widened-Core evidence map.

The review also separates two counting units. Bibliographic screening is
paper-level. Evidence-map coding is row-level, using
\texttt{system\ /\ environment\ configuration} as the unit of analysis.
This distinction is necessary because one paper or system family may
expose more than one spatial interface. The current screening summary
starts from \texttt{417} screened records, assigning \texttt{12} to
\texttt{Core}, \texttt{42} to \texttt{Adjacent}, \texttt{47} to
\texttt{Foundational}, and \texttt{316} to exclusion categories. After
full-text rechecks, targeted widening, and boundary adjudication, the
stable widened Core used for the evidence map contains \texttt{32}
paper-level sources and \texttt{34} coded rows, with \texttt{19}
\texttt{anchor\_core} rows and \texttt{15} \texttt{bridge\_core} rows.

The closure status is also part of the scope boundary. Most rows now
have local PDF or full-text-derived closure cards, but \texttt{BK02}
remains a source-note-only bridge row pending full-text acquisition or
downgrade. It is retained only as lower-weight bridge context and should
not carry strong evidence claims.

The representation taxonomy runs from \texttt{L0} to \texttt{L5}.
\texttt{L0} means no spatial input. \texttt{L1} means place or action
labels. \texttt{L2} means semantic scene description without explicit
topology. \texttt{L3} means local relational structure such as
adjacency, co-presence, nearby agents, or local movement options.
\texttt{L4} means agent-facing global abstract structure such as
integration, depth, control, choice, network position, or other
configuration-level information. \texttt{L5} means geometry-bearing or
embodied input consumed by the agent, such as coordinates, visual
fields, physical constraints, or equivalent position-bearing state. The
levels are not a quality ladder. They are a coding scheme for structural
explicitness at the agent interface.

\subsection{1.4 Main Findings Preview}\label{main-findings-preview}

The evidence map supports a careful descriptive finding. Current
LLM-agent social simulations are not spatially empty. Many systems stage
activity in places, neighborhoods, game worlds, social VR environments,
towns, grids, graphs, or 3D scenes. Some report spatially relevant
outcomes. Across the stable widened Core, the evidence map contains
\texttt{19} \texttt{observed\_effect} rows and \texttt{15}
\texttt{designed\_affordance\_only} rows.

The sharper finding is representational. The corpus is concentrated at
local and mid-structure interfaces rather than at explicit
configurational input. Across \texttt{34} coded rows, the representation
distribution is \texttt{L1\ =\ 1}, \texttt{L2\ =\ 8},
\texttt{L3\ =\ 18}, \texttt{L4\ =\ 1}, and \texttt{L5\ =\ 6}. In the
strict \texttt{anchor\_core}, \texttt{15} of \texttt{19} rows are
\texttt{L3}, with no \texttt{L4} rows. Under the widened reading, the
only \texttt{L4} case appears in the bridge layer as a digital-network
case, not as strict anchor evidence for physical-layout Space Syntax
mediation. The correct conclusion is not that space is absent, and not
that the \texttt{L4} gap is solved. The correct conclusion is that
agent-facing global abstract structure remains highly underexplored.

This matters because social behavior claims depend on the representation
level. A local co-presence or adjacency interface can support claims
about local opportunity, encounter, or situated interaction. It does not
automatically support claims about depth, integration, choice, control,
or configuration-wide position. A 3D backend may support embodied
interaction, but it does not automatically mean the LLM receives
geometry. The bridge \texttt{L5} rows are also heterogeneous: platform
support, embodied cooperation, and VR guidance should not be collapsed
into a single social-mediation result. A graph backend may support
analysis, but it does not count as \texttt{L4} unless global abstract
structure is agent-facing.

The review's agenda follows from this mismatch. Adjacent evidence
suggests that richer spatial input is technically plausible, but
feasibility is not validation. Space Syntax offers transferable
hypotheses, but physical-space findings are not direct evidence for
LLM-agent societies. Current social-simulation systems show many ways
space can scaffold interaction, but stronger claims require matched
controls. Future studies need to specify what agents receive, compare
representation levels, measure behavior at movement, encounter,
interaction, and macro-social scales, and replicate across layouts,
tasks, models, populations, and seeds.

\subsection{1.5 Contributions and
Organization}\label{contributions-and-organization}

The paper makes three contributions. First, it provides a structured
evidence map of agent-accessible spatial representation in LLM-agent
social-simulation systems, explicitly separating backend richness from
what agents consume. Second, it connects Space Syntax and
spatial-simulation concepts to LLM-agent research as transferable
hypotheses rather than as direct evidence. Third, it proposes a research
agenda and evaluation ladder for moving from spatial affordance to
spatial sensitivity, spatial mediation, and eventually replicated
mechanism.

The rest of the paper is organized accordingly. Section 2 introduces the
minimal Space Syntax concepts needed for the review and states the
transfer boundary. Section 3 presents the review protocol,
representation taxonomy, and main evidence map. Section 4 asks whether
richer spatial input is technically plausible. Section 5 synthesizes how
current LLM-agent social simulations use space. Section 6 defines
evaluation requirements for spatially mediated behavior. Section 7 turns
those requirements into a research agenda. Section 8 concludes with the
main findings, limitations, and implications.

\section{Space Syntax Primer}\label{space-syntax-primer}

This section provides only the Space Syntax concepts needed for the rest
of the survey. Its role is not to reteach the whole Space Syntax
literature, but to give AI and multi-agent-systems readers a precise
vocabulary for configuration, movement opportunity, encounter structure,
and claim boundaries. The section should be read as a theoretical
bridge: physical-space findings motivate hypotheses for LLM-agent
systems, but they are not direct evidence about current LLM-agent social
behavior.

\begin{table*}[t]
\centering
\caption{Table 5. Space Syntax Measure Primer for This Survey}
\scriptsize
\setlength{\tabcolsep}{3pt}
\renewcommand{\arraystretch}{1.12}
\begin{tabular}{@{}>{\raggedright\arraybackslash}p{0.175\textwidth}|>{\raggedright\arraybackslash}p{0.175\textwidth}|>{\raggedright\arraybackslash}p{0.175\textwidth}|>{\raggedright\arraybackslash}p{0.175\textwidth}|>{\raggedright\arraybackslash}p{0.175\textwidth}@{}}
\toprule
measure & Core intuition & What an LLM-agent system would need to expose & Survey use & Claim boundary \\
\midrule
integration & How accessible a location is relative to the whole layout. Highly integrated locations tend to be easier to reach from many other locations. & A layout-wide accessibility score, rank, or equivalent global position cue available to the agent, or geometry from which this relation can be derived under a controlled design. & Frames hypotheses about likely movement concentration, encounter opportunity, and publicness. & Foundational theory motivates future tests; it is not direct evidence that current LLM-agent societies are integration-mediated. \\
depth & How many steps or transitions separate one location from others. Deeper locations are more segregated from the wider layout. & A global or task-relevant depth / segregation cue, not only a local place label or immediate neighbor list. & Frames hypotheses about privacy, withdrawal, selective interaction, and low-contact behavior. & Current L3 local adjacency is not enough to support a depth-based claim unless the broader structure is agent-facing. \\
control & How much a location mediates access from its immediate neighbors to the rest of the local structure. & Bottleneck, gatekeeping, or local mediation information exposed as structured input, rather than only an analyst-side graph calculation. & Frames hypotheses about monitoring, guarding, brokerage, and constrained route choice. & Researcher-side control metrics do not count as L4 unless agents receive or can use them. \\
choice & How often a location lies on plausible paths between other locations; related to path betweenness. & Route alternatives, path-betweenness, or global route-structure cues available to the agent, or embodied navigation where path alternatives can be controlled. & Frames hypotheses about route allocation, incidental co-presence, and flow concentration. & Mobility in a rich environment is not by itself evidence of choice-based social mediation. \\
visibility / openness & What can be seen or perceived from a location, including line-of-sight and field-of-view constraints. & Visual fields, observable zones, occlusion, openness, or embodied perception supplied to the agent. & Frames hypotheses about awareness, approach, avoidance, coordination, and interruption. & Visibility can be L5 or structured L4/L3 depending on the interface; the level depends on what the agent consumes. \\
\bottomrule
\end{tabular}

\label{tab:table-5-space-syntax-measure-primer-for-this-survey}
\end{table*}

\begin{figure*}[t]
\centering
\includegraphics[width=0.94\textwidth,height=0.62\textheight,keepaspectratio]{../figures/gpt\_image\_2/figure\_5\_local\_global\_claim\_boundary\_gpt\_image\_2\_v3.png}
\caption{Figure 5. Toy example showing why local adjacency is not equivalent to global configuration. If an agent only receives immediate-neighbor information, both focal locations can look structurally similar. Configurational claims require additional agent-facing information about whole-layout position, such as depth, integration, control, or choice. The example motivates the survey's separation between `L3` local relations and `L4` global abstract structure and clarifies why Space Syntax propositions remain hypotheses until global structure is exposed to agents.}
\label{fig:5}
\end{figure*}

\subsection{2.1 Why Configuration Matters
Here}\label{why-configuration-matters-here}

The key idea this survey borrows from Space Syntax is that spatial
behavior is shaped not only by individual places, but by the way places
are configured in relation to the whole layout. A room, street,
corridor, plaza, or virtual location does not have only local
properties. It also has a position in a wider system of possible
movement, visibility, access, and encounter. That whole-layout position
can affect who passes through, who meets whom, where activity
concentrates, and where withdrawal or privacy becomes easier.

This is directly relevant to LLM-agent social simulation because many
current systems place agents in environments but expose only a partial
spatial interface to them. A system may contain a map, a 3D world, or a
graph backend while giving the agent only a place name, a scene
description, a nearby-agent list, or a set of local movement options.
Space Syntax helps articulate what is missing from such interfaces:
configuration-wide structure that summarizes how a location sits within
the broader environment.

The transfer boundary is essential. Space Syntax findings in physical
environments can motivate hypotheses about movement, co-presence,
encounter probability, access asymmetry, privacy, and route choice. They
do not by themselves show that LLM-agent societies are already
configuration-mediated. For that claim to be testable, the relevant
structure must be exposed to the agent and linked to reported behavior
under an appropriate comparison.

\subsection{2.2 Measures the Reader
Needs}\label{measures-the-reader-needs}

The rest of the survey uses a small set of Space Syntax concepts.
\texttt{Integration} describes how accessible a location is relative to
the whole layout. Highly integrated locations are easier to reach from
many other locations and are often hypothesized to support movement
concentration, publicness, and encounter opportunity. In an LLM-agent
system, an integration-based claim would require agents to receive a
layout-wide accessibility cue, a derived rank, or geometry from which
such a relation can be meaningfully used under controlled conditions.

\texttt{Depth} describes how many steps or transitions separate one
location from others. A deeper location is more segregated from the rest
of the layout. In the context of LLM-agent systems, depth-related
hypotheses would concern privacy, withdrawal, selective interaction, or
reduced incidental contact. A local place label or immediate neighbor
list is not enough to support a depth-based claim unless the broader
structure is agent-facing.

\texttt{Control} describes how much a location mediates access between
its neighbors and the surrounding structure. In social terms, control
can motivate hypotheses about monitoring, guarding, brokerage,
chokepoints, and constrained route choice. For this survey, the
important point is not the formula but the interface requirement: a
researcher-side control calculation does not count as agent-accessible
\texttt{L4} unless the agent receives or can use that mediation
information.

\texttt{Choice} describes how often a location lies on plausible paths
between other locations; it is related to path betweenness. Choice is
useful for thinking about route allocation, incidental co-presence, and
flow concentration. Again, mobility in a rich environment is not itself
evidence of choice-based social mediation. The agent must receive route
alternatives, path-structure cues, or embodied geometry in a design
where those cues can be isolated.

Finally, visibility or openness describes what can be seen or perceived
from a location, including line-of-sight, field-of-view, occlusion, or
observable zones. This concept may enter an LLM-agent system as
\texttt{L5} visual or embodied input, or as a more symbolic structured
cue, depending on what the agent consumes. The level is not determined
by the existence of a visual engine. It is determined by the
agent-facing interface.

Table 5 summarizes these concepts in the terms used by this survey. The
table deliberately includes a claim-boundary column because each measure
can easily be overclaimed. The measures are useful because they define
testable representational targets. They do not convert physical-space
evidence into direct evidence about LLM-agent social behavior.

\subsection{2.3 Local Adjacency Is Not Global
Configuration}\label{local-adjacency-is-not-global-configuration}

The most important distinction for the evidence map is the difference
between local adjacency and global configuration. Two locations can
expose the same local relation to an agent while occupying different
positions in the whole environment. For example, both locations may have
two immediate neighbors. If the agent receives only those immediate
neighbors, the two locations can look equivalent from the agent's
perspective. Yet one may be near a main spine, shallow in the broader
layout, or positioned along many likely routes, while the other may sit
deeper in a side branch.

This distinction is why the survey separates \texttt{L3} from
\texttt{L4}. \texttt{L3} captures local relations: adjacency,
co-presence, nearby agents, local movement options, and local graph
exposure. These are important because they structure immediate
interaction opportunities. But \texttt{L3} does not automatically
provide whole-layout position. Claims about depth, integration, control,
or choice require additional information about the broader
configuration.

Figure 5 visualizes this claim boundary. Panel A shows local views that
can look equivalent. Panel B reveals their different whole-layout
positions. Panel C states the resulting inference rule: local
opportunity or co-presence claims can often be discussed at \texttt{L3},
but configuration-wide claims require \texttt{L4} or controlled
geometry-bearing \texttt{L5}. This is an explanatory diagram, not
evidence from a coded system.

The implication is methodological. If a paper reports that agents move,
meet, or coordinate in a spatial environment, the review still asks what
spatial structure the agent received. A global layout may exist in the
simulator. A researcher may compute graph metrics after the fact.
Neither condition is enough for a configurational agent-facing claim
unless the relevant structure was part of the agent's decision
interface.

\subsection{2.4 Transfer Boundary for the
Survey}\label{transfer-boundary-for-the-survey}

Space Syntax enters this survey as a source of transferable propositions
and missing representation layers. It supplies a vocabulary for asking
whether accessibility, segregation, mediation, route structure, or
visibility could affect LLM-agent movement and interaction if such
structure were exposed to agents. It does not supply direct validation
of current LLM-agent systems.

This boundary shapes the rest of the paper. In Section 3, the evidence
map asks whether current systems expose the relevant structures. The
answer is mostly no for \texttt{L4}: agent-facing global abstract
structure is absent from the strict anchor core and appears only once in
the stable widened Core, as a digital-network bridge case rather than a
physical-layout validation case. In Section 5, Space Syntax propositions
are therefore treated as hypotheses for future social-simulation
research, not as conclusions already supported by the LLM-agent corpus.
In Section 6, the same distinction becomes an evaluation requirement:
stronger spatial-behavior claims require matched controls over what
agents receive.

The practical rule is the same throughout the survey. If a system
exposes only labels, semantic scene descriptions, or local co-presence,
then configuration-level claims remain untested no matter how rich the
backend world may be. If a system exposes global abstract structure,
geometry, or embodied perception to agents, then stronger spatial
hypotheses become testable, but still require observed behavioral
evidence and appropriate controls before mechanism language is
justified.

\section{Evidence Map}\label{evidence-map}

This section maps how current LLM-agent systems represent space at the
agent interface, what kinds of social behavior they study, and what kind
of evidence they actually provide. Its purpose is descriptive and
gap-oriented. It does not argue that configurational spatial structure
has already been robustly shown to shape LLM-agent social behavior.
Instead, it establishes the current coverage, the current absences, and
the evidential limits that later sections use to motivate feasibility
analysis, social-simulation synthesis, and future evaluation
requirements.

\begin{figure*}[t]
\centering
\includegraphics[width=0.94\textwidth,height=0.62\textheight,keepaspectratio]{../figures/gpt\_image\_2/figure\_2\_record\_to\_row\_pipeline\_gpt\_image\_2\_v3.png}
\caption{Figure 2. PRISMA-ScR screening and evidence-map stabilization flow. The bibliographic screening stage starts from `417` screened records and separates `Core`, `Adjacent`, `Foundational`, and excluded records. The evidence-map stage then operates at the `system / environment configuration` level, preserving the strict `anchor\_core` baseline while adding selected `bridge\_core` rows under widened scope rules. Counts therefore differ by layer: the final stable widened Core contains `32` paper-level sources and `34` coded rows. `HC01` and `TW-02` are retained as boundary materials rather than counted in the stable widened-Core evidence map; `BK02` remains a source-note-only bridge row until resolved.}
\label{fig:2}
\end{figure*}

\begin{table*}[t]
\centering
\caption{Table 2. Review Protocol Summary}
\scriptsize
\setlength{\tabcolsep}{3pt}
\renewcommand{\arraystretch}{1.12}
\begin{tabular}{@{}>{\raggedright\arraybackslash}p{0.30\textwidth}|>{\raggedright\arraybackslash}p{0.30\textwidth}|>{\raggedright\arraybackslash}p{0.30\textwidth}@{}}
\toprule
Protocol element & Main-text summary & Appendix / source \\
\midrule
Review type & Scoping review and structured evidence map, not an effect-size synthesis. & paper/appendix/review\_protocol.md \\
Review aim & Map how LLM-agent systems expose space at the agent interface, what social-behavior scales they study, and what evidence status their results can support. & paper/appendix/review\_protocol.md \\
Bibliographic screening & Phase 1 screening separates records into Core, Adjacent, Foundational, and Excluded; PRISMA-ScR reporting uses bibliographic counts. & results/logs/prisma\_summary.json \\
Corpus roles & Core supports the main evidence map; Adjacent supports feasibility and boundary discussion; Foundational supports theory and transferable hypotheses. & paper/appendix/review\_protocol.md; docs/plans/claim\_matrix.md \\
Core layers & Stable widened Core distinguishes anchor\_core from bridge\_core; bridge rows extend coverage but do not carry the same evidential weight as the strict anchor nucleus. & paper/appendix/review\_protocol.md; docs/plans/claim\_matrix.md \\
Unit of analysis & Screening is paper-level; evidence-map coding is at the system / environment configuration level. Split-row families such as Concordia and SimWorld explain why 32 paper-level sources produce 34 coded rows. & paper/appendix/review\_protocol.md; paper/appendix/appendix\_evidence\_table.csv \\
Spatial coding rule & Code what the agent can consume, not what the simulator stores, renders, or what the analyst computes after the fact. & paper/appendix/review\_protocol.md; paper/appendix/taxonomy\_change\_log.md \\
Representation taxonomy & L0 no spatial input; L1 labels; L2 semantic scene descriptions; L3 local relations; L4 global abstract structure; L5 geometry or embodiment consumed by the agent. & paper/figures/figure\_3\_l0\_l5\_taxonomy.svg; paper/appendix/review\_protocol.md \\
Evidence status & designed\_affordance\_only supports architecture/affordance claims; observed\_effect supports limited reported-association claims; neither alone establishes strong causal mechanism. & paper/appendix/review\_protocol.md; docs/plans/claim\_matrix.md \\
Boundary handling & HC01 is Adjacent / boundary / feasibility evidence. TW-02 is a scope-boundary comparison. Neither enters the stable widened-Core evidence map. BK02 remains source-note-only bridge evidence until full text is acquired or the row is downgraded. & paper/appendix/adjudication\_memo.md; assets/survey\_paper/phase1/phase1\_tw02\_scope\_decision\_2026-04-28.md; assets/survey\_paper/evidence\_closure/global\_consistency\_check\_2026-05-01.md \\
Current stable baseline & Strict anchor baseline: 17 paper-level sources and 19 rows. Stable widened Core: 32 paper-level sources and 34 rows, with 19 anchor\_core rows and 15 bridge\_core rows. The 2026-05-01 closure check leaves only BK02 unresolved as source-note-only bridge evidence. & paper/appendix/appendix\_evidence\_table.csv; paper/tables/table\_3\_core\_evidence\_map.md; assets/survey\_paper/evidence\_closure/global\_consistency\_check\_2026-05-01.md \\
\bottomrule
\end{tabular}

\label{tab:table-2-review-protocol-summary}
\end{table*}

\begin{figure*}[t]
\centering
\includegraphics[width=0.94\textwidth,height=0.62\textheight,keepaspectratio]{../figures/gpt\_image\_2/figure\_3\_agent\_interface\_coding\_system\_gpt\_image\_2\_v3.png}
\caption{Figure 3. Agent-accessible spatial representation taxonomy used in the evidence map. Levels code what the agent can consume, not what the environment stores or what the analyst computes after the simulation. The stable widened Core is concentrated at `L3`, while `L4` appears only once and only in a widened digital-network bridge case. `L5` indicates geometry-bearing or embodied input, not automatic configurational mediation. This supports the interpretation of configurational or globally abstract agent-facing structure as an underexplored design space rather than a validated field-wide layer.}
\label{fig:3}
\end{figure*}

\begin{table*}[t]
\centering
\caption{Table 3A. Representation Coverage by Core Layer}
\scriptsize
\setlength{\tabcolsep}{3pt}
\renewcommand{\arraystretch}{1.12}
\begin{tabular}{@{}>{\raggedright\arraybackslash}p{0.103\textwidth}|>{\raggedright\arraybackslash}p{0.103\textwidth}|>{\raggedright\arraybackslash}p{0.103\textwidth}|>{\raggedright\arraybackslash}p{0.103\textwidth}|>{\raggedright\arraybackslash}p{0.103\textwidth}|>{\raggedright\arraybackslash}p{0.103\textwidth}|>{\raggedright\arraybackslash}p{0.103\textwidth}|>{\raggedright\arraybackslash}p{0.103\textwidth}@{}}
\toprule
core layer & paper-level sources & coded rows & L1 & L2 & L3 & L4 & L5 \\
\midrule
anchor\_core & 17 & 19 & 1 & 0 & 15 & 0 & 3 \\
bridge\_core & 15 & 15 & 0 & 8 & 3 & 1 & 3 \\
stable widened Core & 32 & 34 & 1 & 8 & 18 & 1 & 6 \\
\bottomrule
\end{tabular}
\vspace{0.25em}\par\footnotesize this table summarizes the 2026-05-01 closure baseline as 34 coded rows drawn from 32 paper-level sources. The row-paper difference comes from split-row families such as Concordia and SimWorld.
\label{tab:table-3a-representation-coverage-by-core-layer}
\end{table*}

\begin{table*}[t]
\centering
\caption{Table 3B. Behavioral Scale and Evidence Status}
\scriptsize
\setlength{\tabcolsep}{3pt}
\renewcommand{\arraystretch}{1.12}
\begin{tabular}{@{}>{\raggedright\arraybackslash}p{0.22\textwidth}|>{\raggedright\arraybackslash}p{0.22\textwidth}|>{\raggedright\arraybackslash}p{0.22\textwidth}|>{\raggedright\arraybackslash}p{0.22\textwidth}@{}}
\toprule
behavioral scale & designed\_affordance\_only & observed\_effect & total rows \\
\midrule
interaction & 7 & 6 & 13 \\
emergent\_social\_structure & 4 & 9 & 13 \\
mixed & 4 & 4 & 8 \\
stable widened Core & 15 & 19 & 34 \\
\bottomrule
\end{tabular}

\label{tab:table-3b-behavioral-scale-and-evidence-status}
\end{table*}

\begin{table*}[t]
\centering
\caption{Table 3C. Reading Notes}
\scriptsize
\setlength{\tabcolsep}{3pt}
\renewcommand{\arraystretch}{1.12}
\begin{tabular}{@{}>{\raggedright\arraybackslash}p{0.45\textwidth}|>{\raggedright\arraybackslash}p{0.45\textwidth}@{}}
\toprule
point & interpretation \\
\midrule
Strict nucleus & The anchor\_core layer is concentrated at L3, with one L1 row and three L5 rows but no admitted L2 or L4 rows. \\
Bridge recovery & The widened bridge layer contributes all L2 rows, the only admitted L4 row, and three additional L5 rows, but these rows remain lower-weight bridge evidence. \\
L4 boundary & The only admitted L4 row is a digital-network bridge\_core case, not strict anchor evidence for physical-layout Space Syntax mediation. \\
L5 boundary & The bridge L5 rows are heterogeneous: platform support, embodied cooperation benchmark, and VR guidance/navigation should not be collapsed into one social-mediation result. \\
Closure caveat & BK02 remains source-note-only bridge evidence pending full-text acquisition or downgrade; it should not carry strong evidence claims. \\
Claim discipline & observed\_effect rows show that the corpus is not only design rhetoric, but the evidence remains heterogeneous and uneven across representation levels. \\
Writing rule & Use this table for row-level evidence-map claims. Keep PRISMA-ScR figures on bibliographic screening counts rather than mixing the two levels. \\
\bottomrule
\end{tabular}

\label{tab:table-3c-reading-notes}
\end{table*}

\begin{figure*}[t]
\centering
\includegraphics[width=0.94\textwidth,height=0.62\textheight,keepaspectratio]{../figures/gpt\_image\_2/figure\_4\_evidence\_map\_matrix\_gpt\_image\_2\_v3.png}
\caption{Figure 4. Evidence-map matrix crossing agent-accessible spatial representation with behavioral scale and evidence status. The current literature is not spatially empty: it contains local, semantic, and embodied spatial interfaces and several reported observed effects. The gap is specifically configurational: `L4` appears only once, only in a widened digital-network bridge row, and remains absent from the strict `anchor\_core`. The marginal bars keep `anchor\_core` and `bridge\_core` separate so bridge recovery is not mistaken for strict anchor evidence.}
\label{fig:4}
\end{figure*}

\begin{table*}[t]
\centering
\caption{Table 4. Environment-Side Richness vs Agent-Accessible Structure}
\scriptsize
\setlength{\tabcolsep}{3pt}
\renewcommand{\arraystretch}{1.12}
\begin{tabular}{@{}>{\raggedright\arraybackslash}p{0.22\textwidth}|>{\raggedright\arraybackslash}p{0.22\textwidth}|>{\raggedright\arraybackslash}p{0.22\textwidth}|>{\raggedright\arraybackslash}p{0.22\textwidth}@{}}
\toprule
row & environment-side representation & agent-accessible representation & why this coding matters \\
\midrule
HC06 Project Sid & 3D\_engine & L5 & Direct coordinate-bearing memories and explicit spawn/location state reach the agent interface, so geometry is genuinely agent-facing. \\
HC10 Real-world community HD social simulation & 3D\_engine & L3 & The GIS/BIM/Unreal stack is rich, but the agent still receives categorical, prompt-mediated state rather than direct geometry. \\
HC11 Environment-aware VR roleplay & 3D\_engine & L2 & Coordinate-like fields are treated as structured text-schema context, not as embodied geometry consumed by the agent. \\
HC12A SimWorld visual-GPS split & 3D\_engine & L5 & The embodied split row includes visual plus GPS-like position-bearing input, so geometry is part of the usable interface. \\
HC12B SimWorld scene-graph split & 3D\_engine & L3 & The same system family also exposes an abstracted scene-graph/layout interface, which stays below direct geometry. \\
HC14 Crowd evacuation disaster & graph\_based & L3 & A GIS-derived road network with road attributes does not become L5 when the agent sees only text summaries, nearby communications, and route context. \\
L4R-01 Network formation among multi-LLMs & graph\_based & L4 & This is the one admitted widened bridge case where global abstract structure is truly agent-facing through node degree, neighbors, and community information. \\
BK06 TongSIM & 3D\_engine & L5 & A platform can expose embodied or geometry-bearing interfaces, but this supports interface feasibility more than direct LLM social-simulation mediation. \\
R3-05 CoELA & 3D\_engine & L5 & Embodied cooperation is relevant to spatial coordination, but it is a benchmark setting rather than population-level social simulation. \\
TW-13 TUMSphere & 3D\_engine & L5 & VR guidance and navigation show object-location and route-state integration, not emergent configurational social mediation. \\
\bottomrule
\end{tabular}

\label{tab:table-4-environment-side-richness-vs-agent-accessible-structure}
\end{table*}

\subsection{3.1 Review Protocol Summary}\label{review-protocol-summary}

We frame the review as a scoping review rather than an effect-validation
study. The evidence map therefore asks which spatial representations
appear in current LLM-agent systems, which social behaviors are studied
alongside them, and which claims the available evidence can safely
support. Full protocol details, coding rules, and adjudication materials
are reported in Appendix A; the main text keeps only the protocol
elements needed to interpret the figures and tables in this section.

The review separates three corpus roles before making claims.
\texttt{Core} papers form the main coded corpus for the evidence map:
they describe LLM multi-agent systems with identifiable spatial
environments and socially relevant behavior. \texttt{Adjacent} papers
provide feasibility and boundary evidence about whether current models
may handle richer spatial inputs, but they do not count as direct
evidence for multi-agent social effects. \texttt{Foundational} papers
provide Space Syntax theory, physical-space empirical findings, and
transferable hypotheses; they support the conceptual bridge developed
later in the paper, not direct empirical claims about LLM-agent
societies.

The screening and coding process also has two counting layers. At the
bibliographic screening stage, records are assigned to review roles. The
current screening summary starts from \texttt{417} screened records and
assigns \texttt{12} to \texttt{Core}, \texttt{42} to \texttt{Adjacent},
\texttt{47} to \texttt{Foundational}, and \texttt{316} to exclusion
categories. These numbers should not be read as the final evidence-map
counts. After full-text rechecks, targeted widened review, and boundary
adjudication, the evidence map operates at the
\texttt{system\ /\ environment\ configuration} level rather than at the
paper level.

This unit shift is necessary because several system families expose more
than one agent-facing spatial interface. For example, split-row
treatment prevents a family such as \texttt{Concordia} or
\texttt{SimWorld} from being coded as if all of its configurations
exposed the same spatial information to the agent. The resulting stable
widened Core contains \texttt{32} paper-level sources but \texttt{34}
coded rows. Of these rows, \texttt{19} belong to the strict
\texttt{anchor\_core} and \texttt{15} belong to \texttt{bridge\_core}.
The \texttt{anchor\_core} is the strict nucleus for the cleanest
descriptive mapping and strict-gap claims. The \texttt{bridge\_core}
extends the map with socially and spatially meaningful bridge cases, but
it remains lower weight than the strict nucleus.

Two boundary decisions are especially important for reading the rest of
the paper. \texttt{HC01} is retained as Adjacent and boundary evidence
rather than counted in the stable widened Core. \texttt{TW-02} is
retained only as a scope-boundary comparison and does not enter the
stable widened-Core evidence map. A third closure qualification concerns
\texttt{BK02}: it remains a source-note-only bridge row until full text
is acquired or the row is downgraded. It can orient the widened design
space, but the manuscript should not lean on it for strong evidence
claims. Figure 2 summarizes this record-to-row pipeline so that the
later representation counts are not confused with the earlier PRISMA-ScR
screening categories.

\subsection{3.2 Spatial Representation Taxonomy:
L0-L5}\label{spatial-representation-taxonomy-l0-l5}

The evidence map uses an agent-facing taxonomy of spatial
representation. The central coding rule is simple: code what the agent
can consume, not what the simulator stores or what the analyst computes
after the simulation. This distinction is necessary because visually or
technically rich environments often expose only sparse, local, or
prompt-mediated spatial information at the agent interface.

The taxonomy spans six levels. \texttt{L0} indicates no spatial
information. \texttt{L1} indicates place labels or action-space labels
without explicit spatial relations. \texttt{L2} indicates semantic or
descriptive place information without explicit topology, such as scene
context, named objects, or role-play setting. \texttt{L3} captures local
relational structure, including adjacency, co-presence, nearby agents,
available movement options, local feeds, or local graph exposure.
\texttt{L4} is reserved for agent-facing global abstract structure
beyond local next-step relations, including configurational or
network-position information such as integration, depth, control,
choice, degree, neighborhood, or community information when these are
actually available to the agent. \texttt{L5} requires geometry-bearing
or embodied input consumed by the agent, such as coordinates, visual
field, first-person sensory input, physical constraints, or explicit
position-bearing state.

The taxonomy is not a maturity ladder. Higher levels are not
automatically better for every research question, and richer backends do
not automatically imply richer agent-facing input. A \texttt{3D\_engine}
can support \texttt{L5} if the agent receives coordinates, visual input,
or embodied constraints, but the same backend can remain \texttt{L3} if
the agent receives only categorical local observations or text
summaries. Similarly, a \texttt{graph\_based} environment does not
become \texttt{L4} merely because researchers later compute network
metrics. It reaches \texttt{L4} only when global abstract structure is
part of the agent's decision interface.

This agent-interface rule is the basis for Figure 3 and for the coding
decisions in Table 3. It also explains why the same evidence map can
contain rich virtual environments, graph-based platforms, and
text-mediated worlds without treating them as equivalent. The survey is
not measuring visual fidelity or simulator complexity. It is measuring
the structural explicitness of spatial information that can plausibly
enter the agent's reasoning and behavior generation.

\subsection{3.3 Main Evidence Map}\label{main-evidence-map}

The stable widened-Core evidence map shows a literature concentrated in
local and mid-structure spatial interfaces rather than in explicit
configurational input. Across \texttt{34} coded rows, the representation
distribution is \texttt{L1\ =\ 1}, \texttt{L2\ =\ 8},
\texttt{L3\ =\ 18}, \texttt{L4\ =\ 1}, and \texttt{L5\ =\ 6}. This
distribution makes \texttt{L3} the dominant level: more than half of the
coded rows expose some form of local relational spatial structure, such
as co-presence, adjacency, local movement options, or local graph
relations.

The strict and widened readings must be kept separate. In the
\texttt{anchor\_core}, \texttt{15} of \texttt{19} rows are coded as
\texttt{L3}, with one \texttt{L1} row and three \texttt{L5} rows. No
\texttt{anchor\_core} row reaches \texttt{L2} or \texttt{L4}. The
widened bridge layer changes the descriptive landscape by adding all
\texttt{8} \texttt{L2} rows, three additional \texttt{L5} rows, three
additional \texttt{L3} rows, and the only admitted \texttt{L4} row. This
recovery is important because it shows that the broader design space is
more varied than the strict nucleus alone would suggest. At the same
time, it should not be written as if the bridge layer eliminates the
strict-gap interpretation.

The widened rows also require different qualifiers. The additional
\texttt{L5} bridge rows are useful evidence that geometry-bearing or
embodied interfaces can be connected to LLM or agent systems, but they
are not homogeneous social-simulation validations: \texttt{BK06} is
platform-like, \texttt{R3-05} is an embodied cooperation benchmark, and
\texttt{TW-13} is a VR guidance and navigation system. Similarly, the
single \texttt{L4} row is a digital-network bridge case rather than an
anchor-core spatial-layout case. These rows broaden the map; they do not
by themselves convert the field into one with established
configurational social-mediation evidence.

Behavioral-scale coverage is also uneven. \texttt{Interaction} and
\texttt{emergent\_social\_structure} each account for \texttt{13} rows,
while \texttt{mixed} rows account for \texttt{8}. Evidence status varies
across these scales. Interaction rows include \texttt{7}
\texttt{designed\_affordance\_only} rows and \texttt{6}
\texttt{observed\_effect} rows. Emergent-social-structure rows are more
often reported as observed effects, with \texttt{9}
\texttt{observed\_effect} rows and \texttt{4}
\texttt{designed\_affordance\_only} rows. Mixed rows are evenly split,
with \texttt{4} rows in each evidence-status category. Across the full
stable widened Core, the map contains \texttt{19}
\texttt{observed\_effect} rows and \texttt{15}
\texttt{designed\_affordance\_only} rows.

These observed-effect counts matter, but they do not license strong
causal language. They show that the corpus is not only design rhetoric:
a subset of systems reports spatially relevant behavioral outcomes or
associations. However, the observed-effect rows are distributed unevenly
across representation levels and evidence layers. Many reported effects
occur in local, semantic, embodied, or bridge settings, not in
strict-anchor configurational tests. The evidence map therefore supports
a careful descriptive claim: current LLM-agent systems frequently
include spatial structure and sometimes report spatially relevant
outcomes, but the literature has not yet established configuration-level
social mediation as a robust field-wide result.

Table 3 is the central artifact for this subsection. It should be read
as a row-level evidence map, not as a PRISMA screening table. Its key
contribution is to show where the corpus is dense, where it is thin, and
how claim strength changes when the strict anchor and widened bridge
layers are separated. Figure 4 visualizes the same result as a matrix:
\texttt{L3} is the dense center, \texttt{L4} is nearly empty, and
observed effects exist but remain uneven.

\subsection{3.4 The L4 Gap as an Underexplored Design
Space}\label{the-l4-gap-as-an-underexplored-design-space}

The most important gap in the evidence map is not the absence of space
as such. Space is common in the corpus. The sharper gap is the scarcity
of agent-facing global abstract structure. In the strict
\texttt{anchor\_core}, \texttt{L4} is entirely absent. In the stable
widened Core, \texttt{L4} appears once, in a single admitted
digital-network bridge case: \texttt{L4R-01}, coded as
\texttt{bridge\_core\ /\ L4}.

This single bridge case is analytically useful but should not be
overread. It shows that global abstract agent-facing structure is not
outside the design space. The coded system exposes network-position
information such as node degree, neighbors, common connections,
community labels, or broader network structure to agents in a
decision-relevant format. That is sufficient to prevent the claim that
\texttt{L4} is impossible or conceptually incoherent. It is not
sufficient to claim that the field has systematically operationalized
configurational representation in LLM-agent social simulation, and it
should not be treated as direct evidence that physical-layout Space
Syntax constructs have already been validated in LLM-agent societies.

The strict and widened readings therefore point to the same cautious
conclusion. Strictly, \texttt{L4} is absent from the anchor nucleus.
Under the widened reading, \texttt{L4} appears only at the boundary, in
one digital-network bridge case. The correct interpretation is neither
that there is no \texttt{L4} at all nor that the \texttt{L4} gap has
been solved. The correct interpretation is that configurational or
globally abstract agent-facing representation remains highly
underexplored.

Nearby cases illustrate why the coding rule matters. Some systems
include global network or community metrics in the analysis, but those
metrics remain researcher-side outputs rather than agent-facing inputs.
In such cases, the system can remain \texttt{L3} even when the paper
reports global social-network analysis. The survey is concerned with
what the agent can actually use during behavior generation, not with
what the analyst can compute after observing the simulation.

This \texttt{L4} gap is the bridge to the rest of the paper. Section 4
asks whether current models may be able to process richer
configurational input if it is provided. Section 5 asks how current
social-simulation systems use space despite this gap. Section 6 asks
what future studies would need to report before stronger spatial
mediation claims could be made.

\subsection{3.5 Environment-Side Richness vs Agent-Accessible
Structure}\label{environment-side-richness-vs-agent-accessible-structure}

The evidence map also reveals a recurring mismatch between
environment-side richness and agent-accessible structure. The stable
widened Core contains \texttt{16} rows with \texttt{3D\_engine}
backends, \texttt{9} \texttt{graph\_based} rows, \texttt{5}
\texttt{2D\_grid} rows, and \texttt{4} \texttt{text-only} rows. These
backend categories are useful for describing implementation substrate,
but they do not determine representation level.

Table 4 makes this mismatch concrete. \texttt{Project\ Sid} and the
visual-GPS split of \texttt{SimWorld} are coded as \texttt{L5} because
position-bearing or geometry-bearing information reaches the agent
interface. \texttt{MineLand} similarly reaches \texttt{L5} through
embodied multimodal senses and physical constraints. By contrast, the
real-world community-oriented HD social simulation row remains
\texttt{L3} despite a GIS/BIM/Unreal stack, because the agent-facing
observations remain categorical and prompt-mediated rather than directly
geometric. The scene-graph split of \texttt{SimWorld} also remains below
direct geometry because it exposes an abstracted layout interface rather
than full embodied input. \texttt{HC14} provides a graph-based caution:
a GIS-derived road network with road attributes does not become
\texttt{L5} if the agent sees only textual status summaries, nearby
communications, memories, and route context. The bridge \texttt{L5} rows
add further caution: platform support, embodied cooperation tasks, and
VR navigation are all relevant to interface feasibility, but they should
not be collapsed into a single claim about spatially mediated
multi-agent social emergence.

The one admitted \texttt{L4} bridge case further sharpens the
distinction. \texttt{L4R-01} is graph-based, but it is not coded as
\texttt{L4} simply because a graph exists. It is coded as \texttt{L4}
because global or globally abstract network information is agent-facing
in the decision input. This is the key difference between a graph as
backend substrate, a graph as analyst-side object, and a graph as
agent-accessible representation.

The environment-versus-interface distinction matters for two reasons.
First, it prevents the review from treating simulator richness as
evidence strength. A photorealistic or physically detailed environment
may still provide agents with only low-structure textual cues. Second,
it identifies the missing ingredient in much of the literature. The
problem is not that systems never contain space. The problem is that
many systems do not expose global or geometry-level structure to agents
in a form that can be tied to reported social behavior.

\subsection{3.6 Summary and Transition}\label{summary-and-transition}

The evidence map supports three conclusions. First, current LLM-agent
social simulations are spatially present but representation-limited:
they are concentrated at local and mid-structure interfaces, especially
\texttt{L3}. Second, widened bridge cases broaden the representation
landscape, especially for \texttt{L2} and \texttt{L5}, but they do not
erase the strict-gap interpretation for \texttt{L4}. Third, the corpus
contains reported observed effects, but those effects remain
heterogeneous, unevenly distributed, and insufficient for strong
mechanism claims.

These conclusions define the task for the next sections. The feasibility
question is whether current LLMs can process richer spatial or
configurational information if future systems expose it. The
social-simulation question is how current systems already use space and
where their evidence remains below the configurational layer. The
evaluation question is what controls, baselines, and behavioral measures
would be needed to move from spatial affordance toward spatial
sensitivity, mediation, and eventually stronger replicated mechanism
claims.

\section{Feasibility}\label{feasibility}

This section asks a narrower question than the overall survey question.
It does not ask whether richer spatial representation has already been
shown to shape LLM-agent social behavior. It asks whether current models
and adjacent systems provide enough evidence to treat richer spatial
input, including configurational input, as technically plausible for
future social-simulation research.

The safe conclusion is about \texttt{input-side\ feasibility}, not about
\texttt{behavior-side\ validation}. Adjacent spatial-reasoning
benchmarks, spatially aware agents, embodied systems, and boundary cases
suggest that current models may be able to consume structured spatial
inputs when those inputs are made explicit. They do not show that
LLM-agent societies are already spatially mediated, and they do not
remove the need for the representation controls developed in Sections 6
and 7.

\subsection{4.1 Spatial Reasoning Benchmarks as Feasibility
Evidence}\label{spatial-reasoning-benchmarks-as-feasibility-evidence}

The first source of feasibility evidence comes from \texttt{Adjacent}
work on spatial reasoning benchmarks and spatially aware LLM or VLM
systems. This literature is relevant because it probes whether models
can process topological, relational, route-based, map-based, or
geometry-sensitive inputs. It is not direct evidence for multi-agent
social effects. Its value for this survey is narrower: it makes richer
agent-facing spatial representation a plausible design option rather
than a purely speculative one.

The benchmark story should be read as mixed. Earlier textual benchmarks
such as spatial question answering and multi-hop relation tasks show
that language models can often handle simple spatial relations but
struggle as relational depth, compositionality, or viewpoint complexity
increases. Later work extends this evaluation space toward map-based
questions, qualitative spatial reasoning in simulated rooms,
vision-language spatial grounding, and 3D-aware or geometry-enhanced
models. These lines of work suggest partial structured-input
feasibility, especially when tasks provide stable and legible spatial
descriptions. They also show why feasibility should not be overstated:
success on local or short-horizon spatial tasks does not imply reliable
configurational reasoning over long-horizon social simulations.

For this survey, the important distinction is between spatial language
matching and structured relation maintenance. A model that answers
``left of'' or ``near'' questions may still fail when it must maintain a
chain of relations, compare alternative paths, reason over hidden layout
structure, or infer a global position from local evidence. Conversely,
progress on relation-rich benchmarks supports the idea that future
systems can expose more structured spatial inputs to agents and then
test whether those inputs affect behavior. The correct claim is
therefore conditional: benchmark evidence suggests that structured
spatial input may be technically consumable; it does not establish that
such input already mediates LLM-agent social behavior.

\subsection{4.2 Topological Structure vs Geometry vs
Configuration}\label{topological-structure-vs-geometry-vs-configuration}

Feasibility depends on distinguishing several kinds of spatial input
that are often collapsed together. Semantic scene descriptions, local
topological relations, global abstract configuration, and
geometry-bearing perception are different interface problems. A model
that can interpret a room description is not automatically shown to
reason over layout-wide structure. A model that operates in a 3D
environment is not automatically shown to receive \texttt{L5} geometry.
A system that stores a graph is not automatically shown to expose
\texttt{L4} global abstract structure to the agent.

Semantic and local relational inputs are the most clearly feasible
design primitives in the current corpus. Many systems already expose
place names, scene descriptions, nearby agents, co-presence, adjacency,
or local movement options. These inputs support \texttt{L1} to
\texttt{L3} interfaces and are consistent with the evidence map's
concentration at \texttt{L3}. They are sufficient for many local
interaction questions, but they do not test configuration-wide claims.

Geometry-bearing input is also technically plausible, though less
uniform. Some embodied, game, VR, or navigation-oriented systems expose
coordinates, visual observations, physical constraints, route state, or
spatially conditioned motion. These cases support the feasibility of
\texttt{L5}-style interfaces in limited settings. They should not be
treated as solving the \texttt{L4} problem. Geometry and configuration
are related but not interchangeable: a geometry-rich system may still
hide global abstract structure from the agent, while a graph-based or
text-based system may expose useful configurational summaries without
full embodied perception.

Configurational input remains the least demonstrated category. It refers
to global abstract structure beyond the next local relation:
accessibility, depth, control, choice, network position, community
position, or similar layout-wide summaries. The stable widened-Core map
contains only one \texttt{L4} row, and that row belongs to the bridge
layer rather than the strict anchor core. This makes \texttt{L4} a
plausible but underexplored design target. It is plausible because graph
and spatial-reasoning work suggest that structured inputs can be
surfaced to agents. It is underexplored because current LLM multi-agent
social-simulation systems rarely make such structure agent-facing.

\subsection{4.3 Pattern Matching vs World-Structured
Reasoning}\label{pattern-matching-vs-world-structured-reasoning}

The second caution concerns interpretation. Even when a model performs
well on a structured spatial task, the result does not by itself reveal
whether the model is maintaining a stable spatial model, exploiting
lexical regularities, using memorized spatial priors, following a narrow
prompt pattern, or adapting to a benchmark-specific shortcut. This
uncertainty matters because the survey's downstream question is not only
whether a model can answer spatial questions, but whether a spatial
interface can shape multi-agent behavior under controlled conditions.

The feasibility claim should therefore remain weak and precise. Current
evidence is consistent with partial structured-input feasibility: models
may be able to consume richer spatial inputs if those inputs are
explicit, stable, and task-relevant. It is not yet evidence that models
robustly reason over configuration in a way that supports social
mediation. A system might display correct local behavior without using
the intended spatial structure, or it might use a structural cue only
because the prompt makes the desired behavior obvious.

This is why feasibility and evaluation cannot be separated. If future
studies want to claim spatial sensitivity, they must show that behavior
changes when spatial input changes while non-spatial context remains
controlled. If they want to claim configurational mediation, they must
show that a specific global abstract input changes behavior beyond local
adjacency, semantic labels, or prompt wording. The world-structured
reasoning question therefore becomes an empirical design problem: expose
the structure, ablate it, compare it to weaker representations, and
measure whether behavior changes in a way that survives plausible
confounds.

\subsection{4.4 Existing Spatially Aware Agents and Embodied Bridge
Systems}\label{existing-spatially-aware-agents-and-embodied-bridge-systems}

Beyond benchmarks, the feasibility story is strengthened by adjacent and
bridge systems that already expose spatial context to agents in
interaction, navigation, embodiment, or socially situated environments.
Their importance is not that they resolve the social-effect question.
Their importance is that they show multiple implementation routes by
which spatial structure can be delivered to an LLM or VLM interface.

Spatially aware human-agent work provides one route. Studies such as
spatially aware LLM interaction research ask whether an agent can
recognize and respond to a user's spatial context during conversation.
This is not multi-agent social simulation, and it is not configurational
evidence. It is useful because it shows that spatial context can enter
language-mediated interaction and can be evaluated as part of
interaction quality, presence, responsiveness, or naturalness. That
supports feasibility and motivation, not direct social-effect
validation.

Embodied and VR systems provide another route. Systems such as
SARAH-like spatially aware embodied agents, social VR agents, NPC
systems, and navigation-oriented boundary cases show that user position,
trajectory, visual context, audio context, body motion, route state, or
physical constraints can be connected to agent behavior. Some of these
cases approach \texttt{L5} interfaces; others remain closer to semantic
or local spatial awareness. The coding lesson is the same as in Section
3: a 3D engine, avatar, or VR scene does not automatically imply full
geometry at the agent interface. The question is what the agent receives
and uses.

Bridge-core systems add a third route by showing that socially situated
spatial or network structure can sometimes be exposed in more abstract
form. The single admitted \texttt{L4} bridge case is especially
important because it indicates that global or globally abstract
agent-facing structure is within the design space. However, because it
is a bridge case, it should not be treated as anchor-core validation. It
supports the claim that configurational input is technically plausible
and underexplored; it does not support a broad claim that the field has
already tested configurational mediation.

Together, these systems shift the feasibility question. The issue is no
longer whether it is imaginable to provide richer spatial input to
agents. The issue is how to provide it in a stable, interpretable, and
controlled form, and how to connect it to social behavior without
confusing implementation richness with evidence strength.

\subsection{4.5 Feasibility Assessment for Configurational
Input}\label{feasibility-assessment-for-configurational-input}

The overall assessment is cautious but positive. Current evidence is
consistent with the technical feasibility of richer spatial input. Local
and mid-structure representations are already common enough to count as
existing design primitives. Geometry-bearing interfaces are implemented
in a limited and heterogeneous set of systems. Configurational input
remains the least demonstrated layer, but it is better described as an
open design frontier than as an implausible idea.

This assessment depends on the corpus-role boundary. Adjacent benchmarks
and spatially aware agents support input-side possibility. Bridge-core
systems broaden the design space. Foundational Space Syntax work
motivates which configurational variables may matter. None of these
sources, by itself, proves social-behavior effects in LLM multi-agent
systems. The stable widened-Core evidence map remains the main source
for descriptive claims about current practice, and it shows that
\texttt{L4} is sparse.

The section therefore licenses a transition, not a conclusion. It allows
the survey to move from ``can richer spatial input be provided at all?''
to ``how should it be represented, controlled, and evaluated?'' The
answer is developed in the next sections. Section 5 shows that current
social simulations already use space mostly as a local, semantic, or
interaction scaffold. Section 6 defines what would count as spatial
sensitivity or mediation. Section 7 turns those requirements into a
research agenda focused on representation contracts, matched controls,
multi-level behavior measures, and replication.

The negative boundary is equally important. This section does not
conclude that models understand configuration robustly. It does not
conclude that embodied or 3D systems solve the representation problem.
It does not conclude that richer spatial input already improves social
simulation. Feasibility is a necessary precondition for future work, not
proof of spatially mediated social behavior.

\section{Space in Social Simulation}\label{space-in-social-simulation}

This section synthesizes how space currently functions inside LLM-agent
social-simulation systems. It does not restate the evidence map row by
row. Its purpose is to convert the coded distribution into a readable
account of current practice: current systems do use space, but they
mostly use it through local, semantic, or interaction-level structure
rather than through agent-facing configurational abstraction.

\begin{table*}[t]
\centering
\caption{Table 6. Space Syntax Proposition Transfer Table}
\scriptsize
\setlength{\tabcolsep}{3pt}
\renewcommand{\arraystretch}{1.12}
\begin{tabular}{@{}>{\raggedright\arraybackslash}p{0.175\textwidth}|>{\raggedright\arraybackslash}p{0.175\textwidth}|>{\raggedright\arraybackslash}p{0.175\textwidth}|>{\raggedright\arraybackslash}p{0.175\textwidth}|>{\raggedright\arraybackslash}p{0.175\textwidth}@{}}
\toprule
Space Syntax / spatial-simulation proposition & Transferable LLM-agent hypothesis & Minimum agent-facing representation needed & Current evidence-map status & Allowable manuscript claim \\
\midrule
More integrated locations may increase movement and encounter opportunities. & If agents receive layout-wide accessibility information, they may allocate more movement or interaction toward highly accessible locations. & L4 for explicit global accessibility, or controlled L5 geometry with derived layout conditions. & L4 is absent from strict anchor\_core and appears only once in a widened digital-network bridge case. & Motivates future testing; does not show current LLM-agent configurational mediation. \\
Deeper or more segregated locations may support privacy, reduced encounter, or selective interaction. & If agents receive depth or segregation cues, they may shift sensitive dialogue, withdrawal, or low-contact behavior toward less accessible areas. & L4 depth / segregation cues, or a controlled L3 vs L4 comparison. & Current rows are concentrated at L3; depth-like global cues are not systematically exposed to agents. & Transferable hypothesis only; not current direct evidence. \\
Control points can mediate access between neighboring spaces. & If agents receive control or bottleneck information, they may concentrate guarding, monitoring, routing, or brokerage behavior at high-control positions. & L4 control / bottleneck indicators, or graph input that exposes more than local adjacency. & Current L3 graph/feed rows often expose local relations but usually not agent-facing global control values. & Identifies a missing evaluation target. \\
Choice or path-betweenness can affect route allocation and incidental co-presence. & If agents can use path-choice structure, route decisions and encounter distributions may change under matched layouts. & L4 choice / betweenness structure, or L5 geometry with controlled route alternatives. & Some mobility rows exist, but representation controls are not sufficient to infer configuration-level mediation. & Supports agenda and evaluation design, not settled effect claims. \\
Visibility and openness can shape awareness, approach, avoidance, and coordination. & If agents receive visual fields, openness, or line-of-sight constraints, interaction initiation and coordination may vary by visibility condition. & L5 visual / embodied input, or explicit visibility fields as structured input. & L5 appears in a limited subset of widened-Core rows; bridge L5 rows are heterogeneous and do not by themselves establish social mediation. & Supports feasibility and future controlled tests. \\
Local co-presence and adjacency shape immediate interaction opportunities. & If agents receive nearby-agent or adjacent-place information, interaction choice may be sensitive to local spatial context. & L3 local relations. & L3 is the densest current layer, especially in anchor\_core. & Supports descriptive mapping of current practice; effect claims require reported observed\_effect rows. \\
\bottomrule
\end{tabular}

\label{tab:table-6-space-syntax-proposition-transfer-table}
\end{table*}

\subsection{5.1 From Classical ABM to LLM-Agent
Societies}\label{from-classical-abm-to-llm-agent-societies}

The idea that environments shape social behavior did not begin with LLM
agents. Classical agent-based modeling and adjacent simulation
traditions have long treated environment structure, movement
constraints, local neighborhoods, and encounter opportunities as
behaviorally relevant. Those traditions are useful here because they
show why spatial form can matter for social processes: where agents can
move, who they can encounter, and which positions mediate access can all
affect downstream patterns.

LLM-agent social simulation inherits part of this intuition, but it
changes the problem. The agents are no longer only rule-based entities
moving through an explicitly formalized space. They are language-driven
systems whose behavior is mediated by prompts, memory, social roles,
tool calls, and often a game master, simulator, or environment wrapper.
As a result, the relevant question is not simply whether the environment
contains spatial structure. The relevant question is what part of that
structure reaches the agent at the moment of perception, reasoning,
planning, or action.

This distinction explains why current LLM-agent systems can stage
socially rich worlds without yet demonstrating strong configurational
social effects. A simulated town, campus, city, game world, or online
community may organize agent activity spatially. Yet if the agent sees
only local co-presence, place names, nearby messages, route prompts, or
scene summaries, then the system is using space primarily as a local
interaction scaffold. That can be behaviorally important, but it is not
the same as exposing configuration-wide structure in the sense
introduced in Section 2.

Earlier spatial-simulation traditions therefore provide a comparison
point and a source of transferable questions, not direct evidence about
current LLM-agent effects. They motivate asking whether integration,
depth, control, choice, visibility, or route structure might matter in
language-agent societies. They do not show that current LLM-agent
systems have already tested those mechanisms.

\subsection{5.2 How Space Is Currently
Used}\label{how-space-is-currently-used}

The widened-Core evidence map contains \texttt{34} coded rows.
Behavioral-scale coverage is split across \texttt{interaction} rows,
\texttt{emergent\_social\_structure} rows, and \texttt{mixed} rows, but
the representation side remains concentrated in \texttt{L3}, with
\texttt{18} rows at local-relational level and only one admitted
\texttt{L4} bridge case. One bridge row, \texttt{BK02}, remains
source-note-only pending full-text acquisition or downgrade, so it
should not carry strong claims on its own. Read together, the
distribution suggests that the literature is more mature in staging
social activity in space than in exposing globally structured space to
agents.

The recurring representation patterns are consistent across otherwise
different systems. Some systems use place labels, named zones, rooms, or
action spaces. Others provide semantic scene descriptions, role-play
contexts, nearby objects, or scenario-specific spatial cues. Many expose
local relations: co-presence, nearby agents, adjacency, local movement
options, feeds, follow graphs, or communication neighborhoods. A smaller
subset provides embodied or geometry-bearing input such as visual
observations, coordinates, GPS-like position, or physical constraints.
Very few expose global abstract or configurational structure as
agent-facing input.

This pattern should be read as a substantive finding rather than as a
deficiency of individual systems. Space is not absent from LLM-agent
social simulation. In many cases, it is central to how scenarios are
staged and how interactions become possible. The current limitation is
more specific: space usually enters as local scene structure, local
social proximity, movement context, or embodied affordance. It rarely
enters as an explicit configuration-wide representation that agents can
use to reason about accessibility, depth, control, route choice, or
broader network position.

The distinction also clarifies why visually impressive or physically
rich simulations do not automatically close the gap. A 3D environment
can support embodied interaction, but if the LLM-facing interface
reduces the world to local textual state, the representation remains
below direct geometry or global configuration. Conversely, a graph-based
digital environment can reach \texttt{L4} if global or globally abstract
network information is exposed to the agent. The representation level
follows the agent interface, not the visual substrate.

\subsection{5.3 What the Current Evidence Actually
Supports}\label{what-the-current-evidence-actually-supports}

The coded corpus is not limited to design rhetoric. Across the stable
widened Core, \texttt{19} rows are coded as \texttt{observed\_effect}
and \texttt{15} as \texttt{designed\_affordance\_only}. This means a
subset of current systems reports spatially relevant outcomes,
behavioral differences, interaction patterns, mobility patterns, or
social-structure measures. However, an \texttt{observed\_effect} code is
not the same as a mechanism claim. It indicates that a spatial-behavior
association or outcome was reported in a system-level study; it does not
by itself establish why the effect occurred or whether it generalizes
across systems.

The behavioral-scale distribution sharpens this point.
\texttt{Emergent\_social\_structure} rows are numerous and contain many
observed-effect cases, but they are overwhelmingly concentrated at
\texttt{L3}, with just one \texttt{L4} case and limited \texttt{L5}
coverage. \texttt{Interaction} rows include the bridge-layer \texttt{L2}
recovery and several observed effects, but much of that evidence
concerns spatially situated dialogue, role-play, human-agent
interaction, or local scene context rather than multi-agent
configurational mediation. \texttt{Mixed} rows contribute important
embodied and mobility-oriented cases, including some \texttt{L5}
systems, but those cases still do not by themselves demonstrate that
global spatial structure drives social emergence.

The safest synthesis is therefore two-sided. On one side, current
LLM-agent social simulation has already moved beyond purely non-spatial
dialogue. Many systems organize behavior around places, local movement,
co-presence, proximity, feeds, routes, or embodied environments. Some
report observed spatially relevant outcomes. On the other side, the
available evidence remains heterogeneous and mostly below the
configurational layer. It supports claims about how space currently
enters systems and where limited associations are reported, not claims
that spatial configuration has been robustly shown to shape LLM-agent
societies.

This is why the \texttt{anchor\_core} and \texttt{bridge\_core}
distinction remains important. Bridge cases broaden the representational
landscape, especially for \texttt{L2}, \texttt{L5}, and the single
admitted \texttt{L4} case. But bridge recovery should not be read as
field-wide validation. The additional bridge \texttt{L5} rows are
heterogeneous: one is platform-like, one is an embodied cooperation
benchmark, and one is a VR guidance and navigation system. The single
\texttt{L4} row is a digital-network bridge case rather than strict
anchor evidence for physical-layout configuration. The strict anchor
still shows a corpus concentrated at \texttt{L3}, with no admitted
\texttt{L4} rows. The widened map is more informative than the strict
map alone, but it does not erase the strict-gap conclusion.

\subsection{5.4 Space Syntax Propositions as Transferable
Hypotheses}\label{space-syntax-propositions-as-transferable-hypotheses}

The role of Space Syntax in this section is to organize hypotheses, not
to import conclusions wholesale. Physical-space evidence and earlier
spatial-simulation work suggest that accessibility, segregation, control
points, route choice, and visibility can matter for movement, encounter,
privacy, co-presence, and social organization. In this survey, those
propositions become testable hypotheses for future LLM-agent systems
only when the corresponding spatial structure is exposed to agents.

Table 6 makes this transfer explicit. A proposition about integrated
locations becomes a hypothesis about whether agents who receive
layout-wide accessibility information allocate more movement or
interaction to accessible locations. A proposition about depth becomes a
hypothesis about whether agents with depth or segregation cues shift
sensitive dialogue, withdrawal, or low-contact behavior toward less
accessible areas. A proposition about control becomes a hypothesis about
whether bottleneck or mediation information affects guarding,
monitoring, brokerage, or route choice. In each case, the current
evidence-map status limits what can be claimed now.

The same pattern applies to choice, visibility, and local co-presence.
Choice or path-betweenness can motivate future tests of route allocation
and incidental co-presence, but current mobility evidence is not enough
to infer choice-based mediation without representation controls.
Visibility and openness can motivate tests of awareness, approach,
avoidance, and coordination, but \texttt{L5} input in a limited subset
of systems does not itself establish social mediation. Local co-presence
and adjacency are the strongest current fit because \texttt{L3} is the
densest layer in the evidence map; even there, effect claims require
reported observed-effect rows rather than design affordances alone.

The resulting agenda is not speculative in the weak sense of being
detached from the corpus. It is grounded in a clear mismatch: current
systems often contain enough spatial machinery to make spatial questions
meaningful, but they rarely expose the configuration-wide structures
needed to test stronger Space Syntax-style propositions. Future work can
therefore ask sharper questions than ``does space matter?'' It can ask
which agent-facing spatial structure is exposed, which behavior is
measured, which controls separate semantic context from structural
input, and which claim level the evidence supports. This framing also
prevents the one digital-network \texttt{L4} bridge case from being
overextended into a general Space Syntax validation claim.

The central diagnosis of this section is therefore the same as the
evidence map, stated in social-simulation terms. Current systems already
use space as a scaffold for interaction, mobility, and social
organization. What remains largely missing is direct evidence that
agent-facing global configuration mediates LLM-agent social behavior
under controlled conditions.

\section{Evaluation Dimensions}\label{evaluation-dimensions}

This section asks what would count as credible evidence for spatially
mediated social behavior in future LLM-agent research. It is not a
general evaluation survey. Its narrower role is to translate the gaps
diagnosed in Sections 3 to 5 into concrete evidence requirements: what
representation must be exposed to agents, what behavior must be
measured, what comparison must be made, and what claim the result could
safely support.

\begin{table*}[t]
\centering
\caption{Table 7. Evaluation Dimensions for Spatially Mediated LLM-Agent Behavior}
\scriptsize
\setlength{\tabcolsep}{3pt}
\renewcommand{\arraystretch}{1.12}
\begin{tabular}{@{}>{\raggedright\arraybackslash}p{0.145\textwidth}|>{\raggedright\arraybackslash}p{0.145\textwidth}|>{\raggedright\arraybackslash}p{0.145\textwidth}|>{\raggedright\arraybackslash}p{0.145\textwidth}|>{\raggedright\arraybackslash}p{0.145\textwidth}|>{\raggedright\arraybackslash}p{0.145\textwidth}@{}}
\toprule
Evaluation dimension & Behavior measured & Minimum spatial representation & Required control or comparison & Candidate measures & Strongest claim if passed \\
\midrule
Movement and route choice & Location choice, path selection, dwell time, route diversity & L3 for local movement options; L4 or controlled L5 for configurational claims & Same agents and task under place-label, local-relation, and global/geometry conditions & route distribution, path length, location entropy, destination frequency & Spatial sensitivity; configuration claim only if global or geometry condition is isolated. \\
Co-presence and encounter structure & Who meets whom, where, and how often & L3 for local co-presence; L4 for layout-wide encounter hypotheses & Matched social setting with altered spatial interface but stable personas, goals, and interaction rules & encounter rate, repeated co-presence, location-conditioned contact matrix & Limited observed-effect claim if spatial interface changes encounter outcomes. \\
Interaction allocation across places or network positions & Dialogue initiation, partner choice, cooperation/conflict allocation & L3 for local relation tests; L4 for network-position or configurational allocation tests & Compare local-only input against global-structure input; hold prompts and social affordances constant & interaction counts by place, partner diversity, brokerage behavior, topic/place coupling & Reported association or spatial sensitivity; not mechanism unless controls isolate representation. \\
Privacy, avoidance, and selective disclosure & Sensitive dialogue, withdrawal, avoidance, restricted interaction & L2 for semantic privacy cues; L4 for depth/segregation claims & Separate semantic labels from structural depth or accessibility cues & disclosure rate, avoidance moves, private-location selection, sensitive-topic placement & Transfer from semantic affordance to structural mediation only if semantic and structural cues are separated. \\
Group formation and role differentiation & Clustering, segregation, role emergence, community structure & L3 for local network formation; L4 for global network/layout mediation & Matched runs with different spatial interfaces or layouts; same agent population and task regime & modularity, group persistence, role concentration, cross-group interaction & Limited macro-level observed effect if controls are present; stronger claims need replication across layouts. \\
Embodied coordination and spatial cooperation & Joint tasks, collision avoidance, escorting, navigation-assisted cooperation & L5 & Compare geometry-bearing input against symbolic local relation input; hold task and communication budget fixed & task success, coordination latency, communication cost, spatial errors & Embodied-input effect claim; not automatically configurational mediation or population-level social-simulation evidence. \\
Mechanism and confound separation & Whether behavior changes because of spatial structure rather than prompt semantics or task wording & Depends on target mechanism, usually L3 vs L4 or L3/L4 vs L5 & Matched initial conditions, prompt controls, ablations of labels vs relations vs global structure & effect persistence across seeds, ablation deltas, cross-layout robustness & Mechanism remains tentative unless repeated across systems, models, and layouts. \\
\bottomrule
\end{tabular}

\label{tab:table-7-evaluation-dimensions-for-spatially-mediated-llm-agent-behavior}
\end{table*}

\subsection{6.1 What Would Spatial Behavioral Validity
Mean?}\label{what-would-spatial-behavioral-validity-mean}

For this survey, spatial behavioral validity cannot mean only that a
system has a map, a virtual world, navigation, or place-aware dialogue.
Those are spatial affordances. They show that space is present in the
system design, but they do not by themselves show that agent behavior is
sensitive to spatial structure. A system may let agents move through
places while behavior remains driven mainly by prompts, goals, roles,
scripts, or social context.

The next level is spatial sensitivity. A system shows spatial
sensitivity when behavior changes in response to variation in
agent-facing spatial input. This might involve route choice, dwell time,
co-presence, encounter frequency, dialogue allocation, avoidance,
disclosure, group formation, or role differentiation changing when the
spatial interface changes. Sensitivity is stronger than affordance
because it links spatial input to behavioral output, but it still may
not identify the mechanism.

The strongest target is spatial mediation. A study supports a spatial
mediation claim when it shows that a specific type of agent-facing
spatial structure contributes to a behavioral outcome under matched
controls. For this survey, the important word is specific. A mediation
claim should identify whether the relevant structure is semantic scene
information, local adjacency, global abstract configuration,
geometry-bearing perception, or some controlled combination. It should
also separate that structure from confounds such as richer prompts, more
context, different goals, different tasks, or altered social
affordances.

Most current systems in the evidence map can be discussed safely at the
affordance level. Some \texttt{observed\_effect} rows may support
limited spatial-sensitivity claims in particular settings. Very few, if
any, cleanly establish spatial mediation under matched controls,
especially for \texttt{L4} configurational input. In the current closure
baseline, \texttt{L4} is absent from the strict \texttt{anchor\_core}
and appears only once as a widened digital-network bridge case. That row
is useful for defining a possible evaluation target, but it should not
be treated as validation of physical-layout Space Syntax mediation. This
is not a failure of the field so much as a sign of the next evaluation
problem: future work needs to move from spatial presence toward
controlled spatial-behavior coupling.

\subsection{6.2 Candidate Evaluation
Dimensions}\label{candidate-evaluation-dimensions}

Table 7 organizes candidate evaluation dimensions around behaviors that
matter for spatially mediated social simulation. Each row specifies the
behavior being measured, the minimum spatial representation needed, the
required comparison, candidate measures, and the strongest claim that a
successful result could support. This structure matters because
different behavioral questions require different representation levels.
The table is therefore prospective: it describes what future studies
would need to show, not what the current evidence map already
demonstrates.

Movement and route choice can often begin at \texttt{L3}, because local
movement options may already shape where agents go next. However, a
configurational route claim requires \texttt{L4} or controlled
\texttt{L5}: agents must receive global route structure, path
alternatives, accessibility cues, or embodied geometry in a design that
isolates those inputs. Candidate measures include route distribution,
path length, location entropy, dwell time, and destination frequency.
Passing such a test may support spatial sensitivity; it supports a
configuration claim only if global or geometry-bearing structure is
isolated.

Co-presence and encounter structure are central to the social-simulation
problem. At the local level, \texttt{L3} can support tests of
nearby-agent exposure, adjacency, and contact opportunities. For
layout-wide encounter hypotheses, \texttt{L4} is needed because the
claim concerns how broader structure shapes who meets whom, where, and
how often. Candidate measures include encounter rate, repeated
co-presence, contact matrices conditioned on location, and distribution
of social contact across places or network positions.

Interaction allocation asks whether spatial structure affects who talks,
cooperates, conflicts, or exchanges information with whom. Some versions
of this question can be tested with \texttt{L3} local relation input.
Stronger claims about network-position or configurational allocation
require \texttt{L4}, because the relevant treatment is no longer simply
who is nearby but how agents understand position in a broader structure.
The necessary comparison is local-only input versus global-structure
input while holding prompts, roles, tasks, and social affordances
constant.

Privacy, avoidance, and selective disclosure illustrate why semantic and
structural cues must be separated. An agent may choose a ``private
room'' because a prompt labels it private; that is a semantic affordance
and may be meaningful at \texttt{L2}. A different claim is that agents
shift sensitive dialogue toward deeper or less accessible areas because
they use structural depth or segregation cues. That latter claim
requires \texttt{L4} or a controlled comparison between semantic labels
and structural position.

Group formation, segregation, role differentiation, and macro-level
social structure are the highest-risk evaluation targets. They are
attractive because they connect most directly to emergent social
behavior, but they are also easy to overread. A macro pattern observed
after a simulation is not enough. Future studies would need matched runs
with different spatial interfaces or layouts, stable agent populations,
controlled tasks, and repeated seeds. Candidate measures include
modularity, group persistence, role concentration, cross-group
interaction, and spatial distribution of roles.

Embodied coordination and spatial cooperation form a related but
distinct dimension. These questions often require \texttt{L5}, because
collision avoidance, escorting, joint navigation, and physically
grounded cooperation depend on geometry-bearing input or embodied
perception. A successful comparison can support an embodied-input effect
claim, but it should not automatically be treated as configurational
mediation. This distinction matters for the current bridge \texttt{L5}
rows: platform support, embodied cooperation benchmarks, and VR guidance
systems all show relevant interface possibilities, but they do not
collapse into one generic social-mediation result. Geometry and
configuration are related, but not interchangeable.

\subsection{6.3 Controls, Baselines, and Confound
Separation}\label{controls-baselines-and-confound-separation}

The main methodological challenge is confound separation. If a study
changes spatial representation and unrelated context at the same time,
then the behavioral difference cannot be attributed to the spatial
interface. A stronger evaluation design should hold the agents, goals,
task, social setting, communication budget, and prompt semantics as
stable as possible while varying only the spatial information available
to the agent.

One useful control ladder follows the representation taxonomy. A
\texttt{L1/L2} versus \texttt{L3} comparison asks whether local
relational structure matters beyond place naming or semantic scene
description. A \texttt{L3} versus \texttt{L4} comparison asks whether
global abstract structure matters beyond local adjacency, co-presence,
or local movement options. A \texttt{L3} or \texttt{L4} versus
\texttt{L5} comparison asks whether embodied or geometry-bearing input
changes behavior beyond symbolic structure alone. None of these
comparisons is universally required; the right baseline depends on the
claim being made. But the comparison must match the representation
claim.

The section also cautions against overreading common proxies. Navigation
completion alone does not establish spatial mediation. User satisfaction
alone does not establish spatial mediation. A visually rich environment
does not establish spatial mediation. Emergent macro patterns without
representation controls do not establish spatial mediation.
Source-note-only bridge evidence, including the unresolved \texttt{BK02}
row, also cannot carry strong evaluation claims until full text is
acquired or the row is downgraded. These metrics and materials can be
valuable components of an evaluation package, but they need to be tied
to a controlled spatial interface comparison before stronger claims
become available.

This control logic also protects against prompt confounds. If a
global-structure condition simply gives agents longer prompts, more
semantic hints, or more task guidance, then improved behavior may
reflect extra context rather than spatial structure. Similarly, if a
geometry-bearing condition changes both perception and action
affordances, the study should avoid claiming configurational mediation
unless it isolates which part of the interface produced the effect.

\subsection{6.4 Evidence Ladder for Future
Studies}\label{evidence-ladder-for-future-studies}

The evaluation problem can be summarized as an evidence ladder. The
lowest rung is spatial affordance: the system contains places, routes,
scenes, spatially situated interaction, or embodied action. This is
valuable, but it supports design description rather than effect claims.
The next rung is spatial sensitivity: behavior changes when agent-facing
spatial context changes. This supports limited observed-effect claims
when the outcome is reported clearly and the comparison is
interpretable.

The third rung is spatial mediation: a specific spatial representation
contributes to a behavioral outcome in a way that survives matched
controls. This is the level needed for stronger claims about spatially
mediated social behavior. Even here, mechanism language should remain
cautious unless the study separates plausible confounds and reports
enough evidence to distinguish representation effects from prompt, task,
role, or population effects.

The highest rung is replicated mechanism. This would require repeated,
well-controlled findings across layouts, tasks, populations, models,
seeds, and preferably independent systems. The current literature is not
there yet. The point of the ladder is not to dismiss current work, but
to make claim strength explicit. A study can be useful at the affordance
or sensitivity level without being treated as proof of
configuration-mediated social emergence.

Table 7 therefore contributes more than a list of possible metrics. It
defines how future work could move from
\texttt{designed\_affordance\_only} toward stronger
\texttt{observed\_effect} claims, and eventually toward more robust
mechanism claims. It also sets up the research agenda in Section 7: the
agenda should prioritize missing representation layers, matched
controls, behavior measures tied to social outcomes, and replication
across conditions.

\section{Research Agenda}\label{research-agenda}

\begin{figure*}[t]
\centering
\includegraphics[width=0.94\textwidth,height=0.62\textheight,keepaspectratio]{../figures/gpt\_image\_2/figure\_6\_research\_agenda\_evidence\_ladder\_gpt\_image\_2\_v3.png}
\caption{Figure 6. Research agenda as an evidence ladder derived from the evidence map. The current diagnosis is a spatially active but representation-limited literature: `L3` is dense, `L4` is sparse and bridge-only, and observed effects are heterogeneous. Future work should move from explicit representation contracts to mechanism tests, emergent social measures, generalization checks, and application-specific validation. Each step should preserve the survey's claim discipline: future-work directions are not current evidence unless the system exposes the relevant spatial structure and reports observed spatial-behavior associations.}
\label{fig:6}
\end{figure*}

This section translates the evidence-map gaps into future-work
directions. It does not claim that richer spatial representation has
already been shown to improve LLM-agent social simulation. The agenda
follows from the more limited diagnosis developed in Sections 3 to 6:
current systems frequently contain spatial worlds and often expose local
relations, but agent-facing configurational structure remains sparse,
and stronger claims require controls that separate spatial input from
prompts, tasks, roles, and simulator richness.

The agenda is organized around five linked problems: representation,
mechanism, emergence, generalization, and applications. Each problem
corresponds to a different failure mode in the current literature.
Representation asks what spatial structure agents actually receive.
Mechanism asks whether any observed behavioral difference is caused by
spatial structure rather than by language or task confounds. Emergence
asks how micro-level spatial sensitivity could accumulate into
macro-level social patterns. Generalization asks whether effects survive
across layouts, tasks, models, populations, and seeds. Applications ask
how claims should be bounded when spatially situated LLM agents are used
in games, urban simulation, design evaluation, environmental psychology,
or human-agent interaction.

\subsection{7.1 Representation: Making Configuration
Agent-Facing}\label{representation-making-configuration-agent-facing}

The first agenda item is representational. Future systems need to
specify not only that an environment exists, but what spatial structure
reaches the agent interface. The evidence map shows why this matters. In
the stable widened Core, the dominant representation level is
\texttt{L3}, with \texttt{18} of \texttt{34} coded rows exposing local
relational structure such as adjacency, co-presence, nearby agents,
local movement options, or local graph relations. \texttt{L4} appears
only once, in a widened digital-network bridge case, and is absent from
the strict \texttt{anchor\_core}. The unresolved \texttt{BK02} row also
remains source-note-only and should not be used for strong claims. This
pattern makes configurational input a design frontier rather than an
established practice.

A useful research direction is therefore to develop explicit encodings
for configuration. These encodings do not need to copy Space Syntax
software outputs mechanically. They need to expose global abstract
structure in a form that language agents can use and that researchers
can control. Candidate encodings include ranked accessibility cues,
depth or segregation labels, bottleneck or control indicators,
path-choice summaries, neighborhood or community position, and compact
graph summaries. The existing digital-network \texttt{L4} bridge case
shows that agent-facing global abstract structure is possible, not that
physical-layout configuration has already been validated. For embodied
or geometry-bearing systems, the same agenda may involve controlled
visual fields, coordinates, route alternatives, visibility masks, or
derived layout features.

The central requirement is that the representation be agent-facing and
auditable. If a simulator stores a graph but the agent receives only a
scene description, then the experiment has not tested graph-level
configurational input. If an analyst computes integration or betweenness
after the run, then those measures may explain a pattern, but they were
not part of the agent's decision interface. Future papers should
therefore report a representation contract: the backend spatial
substrate, the agent-facing spatial input, the update frequency of that
input, and any researcher-side metrics computed after the simulation.

This agenda also requires negative and reduced representations. A strong
study should not only test a rich condition. It should include
lower-information baselines such as place labels, semantic descriptions,
local adjacency, or local co-presence. The most informative comparison
is often not ``spatial world versus no spatial world,'' but \texttt{L2}
versus \texttt{L3}, \texttt{L3} versus \texttt{L4}, or symbolic
structure versus geometry-bearing \texttt{L5}. These contrasts allow
future work to identify which spatial layer is doing work.

\subsection{7.2 Mechanism: Separating Spatial Structure from Language
Confounds}\label{mechanism-separating-spatial-structure-from-language-confounds}

The second agenda item is mechanism. The current evidence base contains
\texttt{observed\_effect} rows, but an observed effect is not a
mechanism. It indicates that a spatially relevant outcome or association
was reported; it does not establish that the outcome was caused by
configurational structure. Future research needs designs that can
separate spatial structure from other sources of behavioral change.

The most immediate risk is prompt confounding. A condition that gives
agents global spatial summaries may also give them longer prompts,
richer semantics, clearer task hints, or more social cues. If behavior
changes, the difference may reflect additional context rather than
spatial configuration. A stronger design should hold personas, goals,
tasks, communication rules, memory access, and prompt length as stable
as possible while varying the spatial information itself. Where prompt
length cannot be held constant, studies should report the asymmetry and
include ablations that remove non-spatial hints.

Another risk is semantic-structural conflation. For example, an agent
may choose a ``private office'' because the label says private, not
because the location is deep, segregated, or low in through-movement
potential. Both effects can be valid, but they support different claims.
Semantic scene cues are closer to \texttt{L2}; depth, segregation,
control, and choice are closer to \texttt{L4} when they are
agent-facing. Future studies should therefore separate label effects
from structural effects by crossing semantic labels with structural
positions: a private label in an accessible location, a neutral label in
a deep location, and matched alternatives in between.

Mechanism work also needs stronger ablation logic. A plausible ladder is
to begin with local-relation tests, then add global abstract structure,
then add geometry or embodied perception if the research question
requires it. If \texttt{L3} and \texttt{L4} produce similar outcomes,
then local adjacency or co-presence may be sufficient for the measured
behavior. If \texttt{L4} changes outcomes under matched controls, then a
spatial-sensitivity claim becomes stronger. If the effect persists
across layouts and tasks, then mediation language becomes more
plausible, though still not automatically causal in the strong sense.
The ladder should also keep \texttt{L5} separate from \texttt{L4}:
embodied or geometry-bearing input can change behavior without proving
that agents used global configurational structure.

\subsection{7.3 Emergence: From Micro Spatial Bias to Macro Social
Structure}\label{emergence-from-micro-spatial-bias-to-macro-social-structure}

The third agenda item is emergence. Many of the motivating claims in
this area concern macro-level social outcomes: group formation,
segregation, role differentiation, information diffusion, norms,
brokerage, or persistent interaction networks. These outcomes are
attractive because they connect spatial structure to the
social-simulation promise of LLM agents. They are also easy to
overclaim.

Future studies should make the micro-to-macro pathway explicit. A
configuration may first affect movement options, route choice, dwell
time, or local co-presence. Those micro differences may then alter
encounter rates, partner choice, topic exposure, cooperation
opportunities, conflict frequency, or disclosure patterns. Over repeated
interactions, those interaction differences may accumulate into group
structure, role concentration, unequal access to information, or
persistent social clustering. Without this chain, a macro pattern
observed at the end of a simulation is difficult to interpret.

This suggests a multi-level measurement design. At the movement level,
studies can measure route distribution, location entropy, dwell time,
path length, and destination frequency. At the encounter level, they can
measure contact matrices, repeated co-presence, encounter rate by
location, and exposure diversity. At the interaction level, they can
measure dialogue initiation, partner choice, cooperation, conflict,
brokerage, and topic-place coupling. At the macro level, they can
measure modularity, group persistence, role concentration, cross-group
interaction, or inequality in access to information.

The stronger designs will connect these levels rather than report them
separately. For instance, if a global accessibility cue changes dwell
time in integrated locations, and dwell time changes repeated
co-presence, and repeated co-presence changes group formation, then the
study can make a clearer spatial-sensitivity or tentative mediation
claim. If only the final network modularity changes, the claim remains
weaker because the pathway is underspecified.

The agenda should also include null and boundary results. It may turn
out that some LLM-agent behaviors are less sensitive to spatial
structure than expected because prompts, goals, memories, or social
roles dominate the decision process. Such results would still be
valuable if they are produced by clean representation controls. The
field needs to know not only where spatial configuration matters, but
also when local context, semantic labels, or non-spatial social
instructions override it.

\subsection{7.4 Generalization: Layouts, Models, Populations, and
Seeds}\label{generalization-layouts-models-populations-and-seeds}

The fourth agenda item is generalization. A single successful simulation
is not enough to establish a robust spatial-behavior relationship.
LLM-agent systems are sensitive to model choice, prompt wording, memory
architecture, tool interface, population design, random seed, and
scenario framing. Spatial effects may also depend heavily on the layout
family being tested. Future work should therefore treat replication
across conditions as part of the evidence requirement rather than as a
late-stage robustness check.

Layout generalization is especially important for configurational
claims. A result observed in one graph, town, building, or virtual world
may be a property of that specific layout rather than of the spatial
measure being tested. Studies should vary layouts while preserving the
target contrast: high versus low integration, shallow versus deep
positions, high versus low control, high versus low choice, open versus
occluded visibility, or local-only versus global-structure input.
Synthetic layouts can be useful here because they allow controlled
manipulation, while more realistic layouts can test ecological
plausibility.

Model generalization is equally important. A behavior that appears under
one LLM may disappear under another because models differ in spatial
reasoning, instruction following, memory use, social priors, and
sensitivity to prompt format. Future studies should report model
versions, decoding settings, prompt templates, and memory architecture.
Where possible, they should repeat key conditions across more than one
model family or at least across model sizes. The claim should then match
the evidence: a single-model result supports a local system finding, not
a field-wide conclusion.

Population and task generalization also deserve attention. Spatial
structure may matter differently for cooperative tasks, conflict tasks,
information diffusion, private disclosure, navigation-assisted
coordination, or open-ended role-play. It may also interact with agent
heterogeneity: roles, goals, social ties, mobility preferences, risk
preferences, or prior knowledge. A stronger agenda will test whether
spatial sensitivity persists when the population composition and task
regime change.

Finally, future studies should report seed-level variation. Multi-agent
simulations can produce unstable macro outcomes even when initial
prompts are unchanged. If a claimed spatial effect appears in one run
but not across repeated seeds, it should be described as exploratory. If
it appears across seeds, layouts, models, and tasks, then the claim can
move upward on the evidence ladder toward replicated mechanism.

\subsection{7.5 Applications: Bounding Claims by Representation and
Evidence}\label{applications-bounding-claims-by-representation-and-evidence}

The final agenda item concerns applications. Spatially situated LLM
agents are likely to be used in game NPCs, virtual worlds, urban
simulation, planning support, environmental psychology, architectural
design evaluation, education, training, and human-agent interaction.
These applications motivate the research program, but they also increase
the risk of overclaiming. A visually compelling simulation can look
socially plausible even when the agent-facing spatial input is limited
and the behavioral evidence is weak.

Application papers should therefore state what kind of spatial claim
they are making. A game NPC system may only need local co-presence,
place semantics, and short-horizon movement to improve interaction
quality. That is a valid \texttt{L2} or \texttt{L3} application claim,
not a configurational mediation claim. An embodied cooperation
benchmark, platform demonstration, or VR navigation aid may support an
\texttt{L5} interface or usability claim without proving multi-agent
social mediation. An urban or architectural simulation that aims to
reason about encounter structure, publicness, privacy, or circulation
needs stronger representation and evaluation controls. It should expose
layout-wide or geometry-bearing structure to agents and report behavior
measures tied to the claimed social process.

Urban simulation and environmental psychology applications require
particular caution. Space Syntax and adjacent physical-space studies can
motivate hypotheses about movement, encounter, privacy, visibility, and
accessibility, but those hypotheses cannot be transferred directly into
LLM-agent conclusions. If LLM agents are used as simulated occupants,
citizens, pedestrians, or community members, the study should report
which human behavior theory is being approximated, which agent interface
implements the spatial condition, and which outcomes are being validated
against human, empirical, or design-relevant evidence.

Human-agent and social VR applications raise a related boundary. A
system may show that spatially situated agents are engaging, useful, or
believable to users. That supports an interaction-design or
user-experience claim. It does not by itself show that spatial
configuration mediates multi-agent social behavior. To make the latter
claim, the system would need agent-facing spatial variation, behavioral
measures, and matched comparisons of the kind described in Section 6.

The practical recommendation is to tie application claims to three
labels: representation level, behavioral scale, and evidence status. A
paper might claim, for example, that an \texttt{L3} local-relation
interface supports spatially situated dialogue in a particular NPC
setting, or that an \texttt{L5} embodied interface changes coordination
latency in a cooperative task. Those are clearer and more defensible
claims than saying that a system ``uses space'' or ``models spatial
social behavior'' without specifying what the agents received and what
evidence was observed.

\subsection{7.6 Summary: From Gap to Research
Program}\label{summary-from-gap-to-research-program}

The agenda can be summarized as a controlled progression. First, future
systems should make spatial representation explicit at the agent
interface. Second, they should compare representation levels under
matched controls. Third, they should measure behavior at movement,
encounter, interaction, and macro-social scales. Fourth, they should
replicate across layouts, tasks, populations, models, and seeds. Fifth,
they should bind application claims to the representation and evidence
actually tested.

This progression matches the evidence ladder in Section 6. Spatial
affordance is the starting point: agents inhabit or interact within a
spatial environment. Spatial sensitivity is the next target: behavior
changes when agent-facing spatial input changes. Spatial mediation is
stronger: a specific spatial representation contributes to a behavioral
outcome under matched controls. Replicated mechanism is stronger still:
the relationship persists across conditions and can be explained without
relying on prompt, task, or simulator confounds.

The most important near-term target is not to prove a broad thesis that
space shapes LLM-agent societies. The better target is narrower and more
actionable: test whether specific agent-facing spatial structures,
especially underexplored configurational structures, change specific
social behaviors under controlled conditions. If future work follows
that path, the field can move from spatially rich demonstrations toward
cumulative evidence about when, how, and for which behaviors spatial
configuration matters.

\section{Conclusion}\label{conclusion}

This review asked what is known, what is missing, and what is needed to
study how spatial configuration may shape social behavior in LLM
multi-agent systems. The answer is necessarily cautious. Current systems
often contain spatial environments, and many use space to stage
movement, encounter, dialogue, co-presence, or embodied action. Some
systems report spatially relevant outcomes. However, the current
literature does not yet establish that agent-facing spatial
configuration robustly mediates LLM-agent social behavior. The field is
spatially active, but configuration-level evidence remains sparse.

The first conclusion is methodological: future claims need to
distinguish environment-side richness from agent-accessible spatial
representation. A simulator may contain a map, graph, grid, 3D world, or
complex backend while exposing only a place label, scene description,
local nearby-agent list, or prompt-mediated movement option to the LLM.
Conversely, a comparatively simple graph or text interface may expose
structurally meaningful global information if that information is
actually part of the agent's decision input. This is why the review
codes what the agent receives, not what the system stores, renders, or
computes after the fact.

The second conclusion is empirical and descriptive. In the stable
widened-Core evidence map, representation is concentrated at local and
mid-structure levels. \texttt{L3} local-relational input is the dense
center of the corpus, while \texttt{L4} global abstract or
configurational input is absent from the strict \texttt{anchor\_core}
and appears only once in a widened digital-network bridge case.
\texttt{L5} geometry-bearing and embodied inputs appear in a limited
subset of systems, but they are heterogeneous and do not automatically
solve the configurational-representation problem. This pattern means
that the main gap is not the absence of space. The gap is the scarcity
of agent-facing global structure tied to controlled social-behavior
evidence.

The third conclusion concerns claim strength. The evidence map includes
\texttt{observed\_effect} rows, so the literature should not be
described as purely speculative or purely design-oriented. At the same
time, observed effects remain heterogeneous and do not by themselves
establish mechanism. Many claims remain at the level of spatial
affordance or limited spatial sensitivity. Stronger mediation claims
require matched controls over representation, semantics, prompts, tasks,
agent populations, memory, and simulator affordances. Replicated
mechanism claims require robustness across layouts, tasks, models,
populations, and seeds.

The review has several limitations. It is a scoping review of a young
and fast-moving area, so corpus boundaries and bridge-case decisions are
inherently contestable. Some systems are difficult to code because
papers do not always report the exact agent-facing spatial input. The
widened-Core map improves coverage but introduces lower-weight bridge
evidence that must be kept distinct from the strict anchor nucleus. One
bridge row, \texttt{BK02}, remains source-note-only until full text is
acquired or the row is downgraded. The review also relies on Space
Syntax and physical-space research as a source of transferable
hypotheses, not as direct validation for LLM-agent systems. These
limitations are part of the reason the paper emphasizes claim discipline
rather than broad causal conclusions.

The resulting research program is straightforward. Future work should
make spatial representation explicit at the agent interface, compare
representation levels under matched controls, measure behavior across
movement, encounter, interaction, and macro-social scales, and report
robustness across layouts, tasks, models, populations, and seeds.
Configurational input is a particularly important target because it
remains highly underexplored in the current evidence map. The near-term
goal should not be to prove that space generally shapes LLM-agent
societies. It should be to test whether specific agent-facing spatial
structures change specific social behaviors under controlled conditions.

If that program develops, spatially situated LLM agents can move beyond
impressive worlds and plausible demonstrations toward cumulative
evidence. The contribution of this review is to clarify where that
evidence currently stands, where the gap lies, and what would be
required to turn spatially rich simulations into credible studies of
spatially mediated social behavior.

\end{document}
